% concatenated from main.tex
\documentclass[11pt]{article}

\usepackage[margin=1.1in]{geometry}
\usepackage[T1]{fontenc}
\usepackage[utf8]{inputenc}
\usepackage{lmodern}
\IfFileExists{microtype.sty}{\usepackage{microtype}}{}
\usepackage{amsmath,amssymb,amsthm}
\usepackage{booktabs}
\usepackage[numbers,sort&compress]{natbib}
\usepackage[
  colorlinks=true,
  linkcolor=black,
  citecolor=black,
  urlcolor=black,
  pdftitle={U-centering as subset ANOVA: edge regression and higher-order theory},
  pdfauthor={Xianyang Zhang},
  pdfsubject={U-centering, subset ANOVA, edge regression, and Hoeffding-component pairings},
  pdfkeywords={U-centering, subset ANOVA, edge regression, energy distance, Hoeffding decomposition, U-statistic}
]{hyperref}

\allowdisplaybreaks

\theoremstyle{plain}
\newtheorem{theorem}{Theorem}[section]
\newtheorem{proposition}[theorem]{Proposition}
\newtheorem{lemma}[theorem]{Lemma}
\newtheorem{corollary}[theorem]{Corollary}

\theoremstyle{definition}
\newtheorem{definition}[theorem]{Definition}
\newtheorem{example}[theorem]{Example}

\theoremstyle{remark}
\newtheorem{remark}[theorem]{Remark}

\newcommand{\R}{\mathbb R}
\newcommand{\E}{\mathbb E}
\newcommand{\Pp}{\mathbb P}
\newcommand{\Var}{\operatorname{Var}}
\newcommand{\Cov}{\operatorname{Cov}}
\newcommand{\tr}{\operatorname{tr}}
\newcommand{\rank}{\operatorname{rank}}
\newcommand{\diag}{\operatorname{diag}}
\newcommand{\range}{\operatorname{range}}
\newcommand{\one}{\mathbf 1}
\newcommand{\cI}{\mathcal I}
\newcommand{\cU}{\mathcal U}
\newcommand{\cH}{\mathcal H}
\newcommand{\cL}{\mathcal L}
\newcommand{\cV}{\mathcal V}
\newcommand{\cP}{\mathcal P}
\newcommand{\cS}{\mathcal S}
\newcommand{\inner}[2]{\left\langle #1,#2\right\rangle}
\newcommand{\norm}[1]{\left\lVert #1\right\rVert}
\newcommand{\abs}[1]{\left\lvert #1\right\rvert}
\newcommand{\given}{\,\middle|\,}
\newcommand{\subsets}[2]{\binom{#1}{#2}}
\newcommand{\Ustat}{U\text{-statistic}}
\title{U-centering as subset ANOVA:\\ edge regression and higher-order theory}
\author{Xianyang Zhang\thanks{Email: zhangxiany@stat.tamu.edu}\\[2pt]
\normalsize Department of Statistics, Texas A\&M University}
\date{August 8, 2026}

\begin{document}
\maketitle


\begin{abstract}
The unbiased sample versions of squared distance covariance and the
Hilbert--Schmidt independence criterion (HSIC) are fourth-order U-statistics,
yet U-centering evaluates them from pairwise arrays in $O(n^2)$ operations.
We show that
U-centering is exactly the least-squares residual obtained after fitting
additive endpoint effects to a symmetric hollow array.  This interpretation
explains the zero row sums and the denominator $n(n-3)$ through the residual
degrees of freedom.

The same pairwise residualization also gives useful regression
identities.  After endpoint effects are removed from both arrays, the
U-centered dependence $t$-statistic is the ordinary slope $t$-statistic
obtained by regressing one adjusted array on the other.  In the two-sample
problem, pooling the observations
and using the between-group pair indicator as the predictor shows that
the generalized-energy statistic is twice the fitted slope.  The
common-endpoint and fully interacted regressions give the same slope but use
different residual standard errors.

For $n\ge2r$, we extend the construction to arrays indexed by $r$-subsets.
Higher-order U-centering removes all effects involving fewer than $r$ sample
labels, leaves zero $(r-1)$-way margins, and projects onto a residual space of
dimension $\binom nr-\binom n{r-1}$.  For two symmetric kernels with $r$
arguments, the normalized inner product of the centered arrays is unbiased
for the cross-moment of their $r$th Hoeffding components.  A direct estimator
can involve products spanning as many as $2r$ observations, but subset-margin
inversion or higher-order U-centering evaluates the same quantity in
$O(n^r)$ operations for fixed $r$.  When both arrays are formed from the same
kernel and sample, this becomes a nonnegative unbiased estimator of the
variance of the highest-order Hoeffding component.
\end{abstract}

\begin{center}
\small
\textit{Keywords}: ANOVA decomposition, distance covariance, edge regression,
energy distance, Hilbert--Schmidt independence criterion, Hoeffding
decomposition, U-centering, U-statistic, variance component.
\end{center}

\section{Introduction}
\label{sec:introduction}
Many dependence measures are built from pairwise comparisons.  Distance
covariance uses distances between observations
\citep{SzekelyRizzoBakirov2007}, whereas the Hilbert--Schmidt independence
criterion (HSIC) uses entries of kernel Gram matrices
\citep{GrettonEtAl2008,SejdinovicEtAl2013}.  The usual unbiased estimators of
squared distance covariance and HSIC have a useful but initially surprising
property: although their direct formulas involve as many as four distinct
observations, they can be computed from two $n\times n$ arrays in $O(n^2)$
arithmetic
\citep{SzekelyRizzo2014,SongEtAl2012}.  The same algebra appears in martingale
difference divergence and related conditional-mean dependence measures
\citep{ShaoZhang2014,ZhangYaoShao2018}.  For real-valued observations and ordinary
Euclidean distance, the additional ordering structure permits an
$O(n\log n)$ algorithm \citep{HuoSzekely2016}.  In this paper, we instead
study exact centering and inner-product identities for general subset arrays,
without imposing an ordering structure on the observations.  We use the
following notation.  Write
\[
 [n]=\{1,\ldots,n\},
 \qquad
 \binom{[n]}r=\{I\subseteq[n]:|I|=r\}.
\]
For any positive integer $s$, let $\one_s$ denote the all-ones vector in
$\R^s$.
Let $O_i=(X_i,Y_i)$, $i\in[n]$, be independent copies of $O=(X,Y)$, where
$X$ and $Y$ take values in $\mathcal X$ and $\mathcal Y$, respectively.  If
$I=\{i_1<\cdots<i_r\}$, write
$X_I=(X_{i_1},\ldots,X_{i_r})$, and define $Y_I$ similarly.

For a pairwise statistic, the observed entries are indexed by unordered pairs
of distinct sample labels.  We display them as a symmetric \emph{hollow}
matrix, with the absent diagonal cells set to zero.  Let $A=(A_{ij})$ be such
a symmetric hollow array. For $n\ge3$, ordinary U-centering is
\begin{equation}\label{eq:pairwise-ucenter-intro}
 \widetilde A_{ij}
 =A_{ij}
 -\frac{A_{i\cdot}+A_{j\cdot}}{n-2}
 +\frac{A_{\cdot\cdot}}{(n-1)(n-2)},
 \qquad i\ne j,
 \qquad
 \widetilde A_{ii}=0,
\end{equation}
where
\[
 A_{i\cdot}=\sum_{j\in[n]\setminus\{i\}}A_{ij},
 \qquad
 A_{\cdot\cdot}
 =\sum_{\substack{i,j\in[n]\\i\ne j}}A_{ij}
 =2\sum_{1\le i<j\le n}A_{ij}.
\]
For a second symmetric hollow array $B=(B_{ij})$, let $\widetilde B$ denote
its U-centered version, obtained by applying
\eqref{eq:pairwise-ucenter-intro} with $B$ in place of $A$.  When $n=3$, there
are three pairwise observations and three identifiable
endpoint effects.  The additive endpoint model can therefore fit every hollow
symmetric array exactly, so its residual---and hence its U-centered
array---is zero.  For $n\ge4$, the normalized cross-product is
\begin{equation}\label{eq:intro-ucov}
 \frac{1}{n(n-3)}
 \sum_{i\ne j}\widetilde A_{ij}\widetilde B_{ij}.
\end{equation}
For squared distance covariance, unbiased HSIC, and related pair-kernel
statistics, \eqref{eq:intro-ucov} is a fourth-order U-statistic even though
its matrix implementation is quadratic in $n$.

The row sums and grand sum in \eqref{eq:pairwise-ucenter-intro} have a direct
least-squares interpretation.  Treat $A_{ij}$ as an observation on the edge
$\{i,j\}$ of a complete graph and fit
\[
 A_{ij}=u_i+u_j+\varepsilon_{ij}.
\]
The least-squares residual is exactly \eqref{eq:pairwise-ucenter-intro}.  Its
normal equations state that the residuals incident to every label sum to zero,
which explains the zero row and column sums.  Moreover, the residual space has
dimension $\binom n2-n=n(n-3)/2$.  Since the sum in
\eqref{eq:intro-ucov} runs over ordered pairs and therefore counts each edge
twice, its denominator is twice the residual degrees of freedom, namely
$n(n-3)$.

The same residual projector also gives regression interpretations
of dependence and two-sample statistics.  In both problems, pairwise
quantities are adjusted for endpoint effects before a slope is fitted.  For
dependence, both pair arrays are adjusted before one is regressed on the
other.  For two-sample testing, pairwise dissimilarities are regressed on the
between-group pair indicator.  The common-endpoint model assigns one endpoint
effect to each pooled observation across all pair types.  The fully interacted
regression introduced here allows that effect to differ between within- and
between-group pairs and yields the pooled residual standard error used by
\citet{ChakrabortyZhang2021}.  Section~\ref{sec:edge-regression} also relates
the common-endpoint studentization to the pooled U-centered estimator of
\citet{GaoShao2023}, before we turn to higher-order subsets.

We next extend the same construction to larger subsets.  For $n\ge2r$ and an array
$(A_I:I\in\binom{[n]}r)$, write $\cU_{n,r}$ for the higher-order U-centering
operator: $\cU_{n,r}A$ is the unique least-squares residual obtained after
fitting all effects involving fewer than $r$ labels.  Its explicit formula is
given in Definition~\ref{def:higher-u}.  In particular, its $(r-1)$-way
margins vanish:
  \begin{equation*}
 \sum_{i\notin J}(\cU_{n,r}A)_{J\cup\{i\}}=0,
 \qquad J\in\binom{[n]}{r-1}.
\end{equation*}
The residual degrees of freedom are
\[
 d_{n,r}
 =\binom nr-\binom n{r-1}
 =\binom nr\frac{n-2r+1}{n-r+1}.
\]
At $r=1$ this is ordinary centering; at $r=2$ it is
\eqref{eq:pairwise-ucenter-intro}; and at $r=3$ it removes the grand mean,
one-label effects, and two-label effects from a triple-indexed array.

The finite-sample ANOVA above is deterministic, while its statistical
interpretation follows from the population Hoeffding decomposition.  For
symmetric square-integrable kernels $a:\mathcal X^r\to\R$ and
$b:\mathcal Y^r\to\R$, let $a_r$ and $b_r$ denote their $r$th
Hoeffding components (see Section~\ref{sec:general_hoeffding}).  We study
  \begin{equation*}
 \Theta_r(a,b)
 =\E\!\left[a_r(X_1,\ldots,X_r)b_r(Y_1,\ldots,Y_r)\right].
\end{equation*}
For $r\ge1$, if $X$ and $Y$ are independent, then
$a_r(X_1,\ldots,X_r)$ and $b_r(Y_1,\ldots,Y_r)$ are independent.  Since each
$r$th Hoeffding component has mean zero, $\Theta_r(a,b)=0$.
Different component levels can describe different features of dependence, as
the four-argument matching-kernel example in Section~\ref{sec:example}
illustrates.

A direct unbiased estimator groups products $a(X_I)b(Y_J)$ according to the
overlap size $s=|I\cap J|$ and combines the resulting averages with weights
$(-1)^{r-s}\binom rs$.  The weights remove the Hoeffding components indexed
by $j<r$ and retain only the $r$th component.  For fixed $r$, both subset-margin
inversion and the normalized inner product of higher-order U-centered arrays
evaluate the scalar in $O(n^r)$; the latter route also returns the residual
arrays:
  \begin{equation*}
 \frac{1}{d_{n,r}}
 \sum_{I\in\binom{[n]}r}
 (\cU_{n,r}A)_I(\cU_{n,r}B)_I,
 \qquad A_I=a(X_I),\quad B_I=b(Y_I).
\end{equation*}
The projection representation also yields zero margins and residual degrees
of freedom, and shows that centering either array is sufficient.


Hoeffding decompositions and U-statistics are classical
\citep{Hoeffding1948,Serfling1980}.  Finite-population and exchangeable
extensions are developed by \citet{BloznelisGotze2001} and
\citet{Peccati2004}.  The rank theory of inclusion matrices goes back to
\citet{Gottlieb1966}, while projection methods for functions on fixed-size
subsets are discussed by \citet{Dunkl1978,Filmus2016}.  We use this established
projector and rank theory as the algebraic basis for the statistical
development below.  For $n\ge 2r$, \citet[Theorem~1]{WangLindsay2014} give an
exact unbiased estimator of the full variance of a general order-$r$
U-statistic.  Our equal-kernel specialization instead concerns only the
variance of the highest-order Hoeffding component.  In this case, the
projection identity rewrites the overlap-based estimator as a nonnegative
residual mean square.
Against this background, our main contributions are as follows.
\begin{itemize}
\item We show that empirical U-centering is exactly subset ANOVA
residualization.  This interpretation explains the familiar pairwise formula,
zero margins, and residual degrees of freedom, and yields an explicit formula
for general $r$-subset arrays from observable subset margins.
\item At the pairwise level, we give regression representations for dependence
and two-sample homogeneity statistics.  After endpoint effects are removed,
the U-centered dependence $t$-statistic is the ordinary slope $t$-statistic.
In the two-sample problem, the generalized-energy statistic is twice the slope
on the between-group pair indicator.  The common-endpoint and fully interacted
regressions are nested and estimate the same slope but use different residual
standard errors.
\item Higher-order U-centering turns each kernel array into a zero-margin ANOVA
residual.  The normalized inner product of two such residuals equals the direct
overlap estimator.  For fixed $r$, this scalar can be evaluated in $O(n^r)$
operations either from a shared subset-margin table or by higher-order
U-centering.  The same projection gives a regression formulation for
$r$-subset arrays after all lower-order subset effects have been removed.
\item As an application, we show that the overlap-based unbiased estimator of
the variance of the highest-order Hoeffding component equals the normalized
squared norm of the higher-order U-centered residual, making its
nonnegativity immediate.
\end{itemize}



Sections~\ref{sec:pairwise-anova} and~\ref{sec:pairwise-hoeffding} develop the
pairwise projection and population U-statistic identities.
Section~\ref{sec:edge-regression} gives the pair-array regressions, and
Section~\ref{sec:higher-order-u-centering} extends the construction to
$r$-subset arrays.  Section~\ref{sec:population-pairings} develops the
population pairing and two $O(n^r)$ evaluations of its unbiased estimator,
then specializes the projection identity to the variance of the highest-order
Hoeffding component.  The
two-sample regression proofs are collected in
Appendix~\ref{app:two-sample-nesting}.

\section{Pairwise U-centering as ANOVA residualization}
\label{sec:pairwise-anova}

We first establish the projection interpretation in the pairwise case.
Because unordered pairs of sample labels are the edges of a complete graph,
a pairwise array can be viewed as an edge-indexed response.  U-centering
removes the additive effects associated with the two endpoints of each edge
and retains the interaction orthogonal to those endpoint effects.

\subsection{Symmetric hollow arrays}
A symmetric pairwise array assigns one value to each unordered pair
$\{i,j\}\in\binom{[n]}2$.  We display these values in an $n\times n$ matrix,
although its diagonal cells are not observations: a U-statistic uses distinct
sample labels, so the pairs $\{i,i\}$ are absent.  Setting the diagonal to
zero embeds the edge-indexed array in a symmetric matrix without introducing
self-pairs.  Let
\[
 \cV_{n,2}
 =\left\{A\in\R^{n\times n}:A=A^\mathsf T,\ A_{ii}=0
   \text{ for every }i\right\}.
\]
For $A=(A_{ij})$ and $B=(B_{ij})$ in $\cV_{n,2}$, define their edge inner
product by
\begin{equation}\label{eq:edge-inner}
  \inner{A}{B}_{E}=\sum_{i<j}A_{ij}B_{ij}.
\end{equation}
The associated edge norm is $\norm{A}_E=\inner{A}{A}_E^{1/2}$.
\begin{definition}[Pairwise U-centering]\label{def:pairwise-u}
For $n\ge3$ and $A\in\cV_{n,2}$, define $\widetilde A=\cU_{n,2}A$ by
\begin{equation}\label{eq:pairwise-u}
 \widetilde A_{ij}
 =A_{ij}-\frac{A_{i\cdot}+A_{j\cdot}}{n-2}
       +\frac{A_{\cdot\cdot}}{(n-1)(n-2)},
 \quad i\ne j,
 \qquad
 \widetilde A_{ii}=0.
\end{equation}
\end{definition}

The two endpoint terms remove their fitted additive contributions, and the
final grand-sum term corrects for their shared overall mean.  Because the
diagonal entries vanish, this notation agrees
with the standard U-centering formula.  The zero-row-sum and idempotence
properties were recorded by \citet{SzekelyRizzo2014}.  We next show that both
properties follow from an exact least-squares representation.

\subsection{The exact ANOVA characterization}

Define edge vectorization by
$\operatorname{vec}_E(A):=a_E=(A_{ij}:i<j)\in\R^{\binom{n}{2}}$; for a second array $B$,
define $\operatorname{vec}_E(B):=b_E=(B_{ij}:i<j)$.  Let
$Z\in\R^{\binom{n}{2}\times n}$ be the unsigned
edge-by-vertex incidence matrix,
  \begin{equation*}
 Z_{\{i,j\},k}=\mathbf 1\{k\in\{i,j\}\}.
\end{equation*}
Each row of $Z$ corresponds to one edge and contains exactly two ones in the
columns of its endpoints.  Consider the additive pair model
\begin{equation}\label{eq:pair-additive-model}
  A_{ij}=u_i+u_j+\varepsilon_{ij},\qquad i<j.
\end{equation}
Equivalently, writing $u_i=\mu/2+\alpha_i$ and imposing
$\sum_i\alpha_i=0$ gives the familiar ANOVA form
\begin{equation}\label{eq:pair-anova-model}
  A_{ij}=\mu+\alpha_i+\alpha_j+\varepsilon_{ij}.
\end{equation}
This fixed-array decomposition is related to the round-robin ANOVA underlying
the Social Relations Model, which separates actor, partner, and relationship
effects \citep{WarnerKennyStoto1979,KennyLaVoie1984}; in our setting symmetry
constrains the actor and partner roles to a common endpoint effect, and
$\varepsilon_{ij}$ is a deterministic least-squares residual.

\begin{theorem}[U-centering is pairwise ANOVA residualization]
\label{thm:pair-anova}
Let $n\ge3$ and $A\in\cV_{n,2}$.  The U-centered edge vector is the
least-squares residual from \eqref{eq:pair-additive-model}:
\begin{equation}\label{eq:pair-projection}
  \operatorname{vec}_{E}(\widetilde A)
  =M_{n,2}a_E,
  \qquad
  M_{n,2}=I_{\binom{n}{2}}-Z(Z^\mathsf TZ)^{-1}Z^\mathsf T.
\end{equation}
The least-squares estimates in the parametrization
\eqref{eq:pair-anova-model} are
\begin{equation}\label{eq:pair-ls}
 \widehat\mu=\frac{A_{\cdot\cdot}}{n(n-1)},
 \qquad
 \widehat\alpha_i
 =\frac{A_{i\cdot}-A_{\cdot\cdot}/n}{n-2}.
\end{equation}
Consequently, the fitted residual is exactly \eqref{eq:pairwise-u}.
\end{theorem}

\begin{proof}
We compute the Gram matrix of the incidence design and apply its inverse to
the row-sum vector.

Each edge contains two vertices, and two distinct vertices belong together to
exactly one edge.  Therefore
  \begin{equation*}
 Z^\mathsf TZ=(n-2)I_n+\one_n\one_n^\mathsf T.
\end{equation*}
The matrix is nonsingular for $n\ge3$, with inverse
\begin{equation}\label{eq:ztz-inverse}
 (Z^\mathsf TZ)^{-1}
 =\frac{1}{n-2}I_n
 -\frac{1}{2(n-1)(n-2)}\one_n\one_n^\mathsf T.
\end{equation}
Since $(Z^\mathsf Ta_E)_i=A_{i\cdot}$ and
$\one_n^\mathsf TZ^\mathsf Ta_E=A_{\cdot\cdot}$, the fitted vertex coefficient is
\[
 \widehat u_i
 =\frac{A_{i\cdot}}{n-2}
  -\frac{A_{\cdot\cdot}}{2(n-1)(n-2)}.
\]
Thus
\[
 A_{ij}-\widehat u_i-\widehat u_j
 =A_{ij}-\frac{A_{i\cdot}+A_{j\cdot}}{n-2}
  +\frac{A_{\cdot\cdot}}{(n-1)(n-2)},
\]
which is \eqref{eq:pairwise-u}.  Finally,
$\widehat\mu=2n^{-1}\sum_i\widehat u_i$ and
$\widehat\alpha_i=\widehat u_i-\widehat\mu/2$ give
\eqref{eq:pair-ls}.
\end{proof}


\begin{corollary}
\label{cor:pair-basic}
Let
\begin{align*}
 \cH_{n,2}
 &=\left\{C\in\cV_{n,2}:C\one_n=0\right\},\\
 \cL_{n,2}
 &=\left\{C\in\cV_{n,2}:C_{ij}=u_i+u_j
      \text{ for some }u\in\R^n\right\}.
\end{align*}
Then the edge space is the orthogonal direct sum
  \begin{equation*}
 \cV_{n,2}=\cL_{n,2}\oplus\cH_{n,2},
 \qquad \cL_{n,2}\perp\cH_{n,2},
\end{equation*}
and
\[
 \range(\cU_{n,2})=\cH_{n,2},
 \qquad
 \ker(\cU_{n,2})=\cL_{n,2}.
\]
In particular,
  \begin{equation*}
 \widetilde A\one_n=0,
 \qquad
 \one_n^\mathsf T\widetilde A=0,
 \qquad
 \cU_{n,2}^2=\cU_{n,2}.
\end{equation*}
Moreover, $\cU_{n,2}$ is self-adjoint under the edge inner product.
Equivalently, $\widetilde A=\cU_{n,2}A$ is the unique symmetric hollow array
such that
\[
 \widetilde A\one_n=0
 \qquad\text{and}\qquad
 A-\widetilde A\in\cL_{n,2}.
\]
\end{corollary}

\noindent The proof is given in
Appendix~\ref{app:proof-pair-basic}.


\subsection{Residual degrees of freedom and the normalization}
The edge space has dimension $\binom n2$, whereas the additive fitted space
has dimension $n$.  It follows that
  \begin{equation*}
 d_{n,2}:=\dim(\cH_{n,2})
 =\binom n2-n
 =\frac{n(n-3)}{2}.
\end{equation*}
The same rank was obtained from a vectorized representation of U-centering by
\citet{ZhuZhangYaoShao2020}.  For $n\ge4$, using the edge inner product defined in
\eqref{eq:edge-inner},
\begin{equation}\label{eq:pair-normalized-inner}
 \frac{1}{d_{n,2}}\inner{\widetilde A}{\widetilde B}_E
 =\frac{1}{n(n-3)}\sum_{i\ne j}\widetilde A_{ij}\widetilde B_{ij}.
\end{equation}
When $A=B$, \eqref{eq:pair-normalized-inner} is the residual mean square from
the endpoint ANOVA\@.  Self-adjointness and idempotence further give the
one-sided-centering identity
  \begin{equation*}
 \inner{\widetilde A}{\widetilde B}_E
 =\inner{A}{\widetilde B}_E
 =\inner{\widetilde A}{B}_E.
\end{equation*}
This identity underlies familiar simplifications for unbiased squared distance
covariance and HSIC \citep{SzekelyRizzo2014,SongEtAl2012}.

By Theorem~\ref{thm:pair-anova}, the least-squares representation gives the
following closed-form inner-product formula:
\[
 \inner{\widetilde A}{\widetilde B}_E
 =a_E^\mathsf TM_{n,2}b_E
 =a_E^\mathsf Tb_E
 -(Z^\mathsf Ta_E)^\mathsf T
   (Z^\mathsf TZ)^{-1}(Z^\mathsf Tb_E).
\]
Here $Z^\mathsf Ta_E=(A_{1\cdot},\ldots,A_{n\cdot})^\mathsf T$, and
\eqref{eq:ztz-inverse} gives
\begin{equation}\label{eq:pair-contraction}
\begin{split}
 \frac{1}{n(n-3)}\sum_{i\ne j}\widetilde A_{ij}\widetilde B_{ij}
 =\frac{1}{n(n-3)}\bigg\{&
 \sum_{i\ne j}A_{ij}B_{ij}
 -\frac{2}{n-2}\sum_{i=1}^n A_{i\cdot}B_{i\cdot}\\
 &+\frac{A_{\cdot\cdot}B_{\cdot\cdot}}{(n-1)(n-2)}\bigg\}.
\end{split}
\end{equation}
Thus the computation requires only the two arrays, their row sums, and their
grand sums.  Once the arrays have been formed, all terms in
\eqref{eq:pair-contraction} can be accumulated in $O(n^2)$ arithmetic.



\section{\texorpdfstring{{Pairwise Hoeffding decomposition and fourth-order U-statistic identity}}{Pairwise Hoeffding decomposition and fourth-order U-statistic identity}}
\label{sec:pairwise-hoeffding}

We now turn from the sample-label projection to its population target.

\subsection{Population interaction and overlap moments}

Let $a:\mathcal X^2\to\R$ and $b:\mathcal Y^2\to\R$ be symmetric kernels such
that $\E\{a(X_1,X_2)^2\}<\infty$ and
$\E\{b(Y_1,Y_2)^2\}<\infty$.  The pairwise Hoeffding decomposition of $a$ is
\begin{equation}\label{eq:pair-hoeffding}
 a(x_1,x_2)=a_0+a_1(x_1)+a_1(x_2)+a_2(x_1,x_2),
\end{equation}
where
  \begin{align*}
 a_0&=\E a(X_1,X_2)\\
 a_1(x)&=\E a(x,X_2)-a_0\\
 a_2(x_1,x_2)
 &=a(x_1,x_2)-\E a(x_1,X)-\E a(X,x_2)+a_0,
\end{align*}
with $\E\{a_2(x,X)\}=0$ for almost every $x$.
Define $b_0,b_1,b_2$ in the same way and set
\begin{equation}\label{eq:theta2}
 \Theta_2(a,b)=\E\{a_2(X_1,X_2)b_2(Y_1,Y_2)\}.
\end{equation}
In general, $\Theta_2(a,b)$ is a cross-moment and may have either sign.  To
express \eqref{eq:theta2} in terms of observable raw kernels, define
  \begin{align*}
 \eta_2&=\E\{a(X_1,X_2)b(Y_1,Y_2)\}\\
 \eta_1&=\E\{a(X_1,X_2)b(Y_1,Y_3)\}\\
 \eta_0
 &=\E\{a(X_1,X_2)b(Y_3,Y_4)\}
  =\E a(X_1,X_2)\,\E b(Y_1,Y_2).
\end{align*}
Here the subscript records the number of shared sample labels.

\begin{proposition}
\label{prop:pair-overlap}
Let
$\Theta_0=a_0b_0$ and
$\Theta_1=\E\{a_1(X_1)b_1(Y_1)\}$.  Then
\[
 \eta_0=\Theta_0,\qquad
 \eta_1=\Theta_0+\Theta_1,\qquad
 \eta_2=\Theta_0+2\Theta_1+\Theta_2,
\]
and hence
\begin{equation}\label{eq:theta2-overlap}
 \Theta_2(a,b)=\eta_2-2\eta_1+\eta_0.
\end{equation}
\end{proposition}

\noindent The proof is given in
Appendix~\ref{app:proof-pair-overlap}.

\subsection{The direct estimator and its U-centered form}

For $n\ge4$, set
\[
 A_{ij}=a(X_i,X_j),\qquad B_{ij}=b(Y_i,Y_j),\qquad i\ne j,
\]
and set $A_{ii}=B_{ii}=0$.  With $(n)_q=n!/(n-q)!$, define
  \begin{align*}
 \widehat\eta_2
 &=\frac{1}{(n)_2}\sum_{i\ne j}A_{ij}B_{ij}\\
 \widehat\eta_1
 &=\frac{1}{(n)_3}\sum_{i,j,k\ \mathrm{all\ distinct}}
     A_{ij}B_{ik}\\
 \widehat\eta_0
 &=\frac{1}{(n)_4}\sum_{i,j,k,\ell\ \mathrm{all\ distinct}}
     A_{ij}B_{k\ell}.
\end{align*}
Each term is unbiased for the corresponding $\eta_s$, so
\begin{equation}\label{eq:direct-fourth}
 T_{n,2}^{\mathrm{dir}}
 =\widehat\eta_2-2\widehat\eta_1+\widehat\eta_0
\end{equation}
is unbiased for $\Theta_2(a,b)$.  After symmetrization it is a fourth-order
U-statistic.

\begin{theorem}[Pairwise compression identity]
\label{thm:pair-compression}
For $n\ge4$,
\begin{equation}\label{eq:pair-compression}
 T_{n,2}^{\mathrm{dir}}
 =\frac{1}{d_{n,2}}\sum_{i<j}\widetilde A_{ij}\widetilde B_{ij}
 =\frac{1}{n(n-3)}\sum_{i\ne j}\widetilde A_{ij}\widetilde B_{ij}.
\end{equation}
Thus this fourth-order U-statistic can be evaluated in $O(n^2)$ arithmetic once
the two pairwise arrays have been formed.
\end{theorem}

\begin{proof}
Expanding the row-sum cross-product and separating the terms with identical
second endpoints gives
\[
 \sum_{i=1}^n A_{i\cdot}B_{i\cdot}
 =(n)_2\widehat\eta_2+(n)_3\widehat\eta_1.
\]
For the product of the two grand sums, the two ordered edges either have the
same two labels, share exactly one label, or are disjoint.  When they have the
same two labels, the second edge has two possible orientations; when they
share one label, there are four choices for the positions of the common
endpoint.  Therefore,
\[
 A_{\cdot\cdot}B_{\cdot\cdot}
 =2(n)_2\widehat\eta_2
  +4(n)_3\widehat\eta_1
  +(n)_4\widehat\eta_0.
\]
Substituting these identities into \eqref{eq:pair-contraction} and simplifying
the three coefficients yields
\[
 \frac{1}{n(n-3)}
 \sum_{i\ne j}\widetilde A_{ij}\widetilde B_{ij}
 =\widehat\eta_2-2\widehat\eta_1+\widehat\eta_0,
\]
which is \eqref{eq:direct-fourth}.
\end{proof}

Theorem~\ref{thm:general-compression} below extends the compression identity to
arbitrary $r$-subsets.  The next proposition gives a complementary pointwise
interpretation: conditionally, each U-centered entry recovers a finite-sample
multiple of the population interaction.

\begin{proposition}
\label{prop:pair-conditional}
For a fixed pair $\{i,j\}$ and $n\ge3$,
\begin{equation}\label{eq:pair-conditional}
 \E\left(\widetilde A_{ij}\mid X_i,X_j\right)
 =\frac{n-3}{n-1}a_2(X_i,X_j).
\end{equation}
\end{proposition}

\noindent The proof is given in
Appendix~\ref{app:proof-pair-conditional}.

\subsection{Distance covariance and HSIC}

For distance covariance, take
\[
 a(x,x')=\norm{x-x'},\qquad b(y,y')=\norm{y-y'}.
\]
Then $a_2$ and $b_2$ are the population-centered distance kernels and
$\Theta_2(a,b)=\operatorname{dCov}^2(X,Y)$.  Under the sufficient moment
condition $\E\norm X^2+\E\norm Y^2<\infty$,
\eqref{eq:pair-compression} is the usual unbiased squared sample distance
covariance \citep{SzekelyRizzoBakirov2007,SzekelyRizzo2014}.

For HSIC, let $a=k$ and $b=\ell$ be positive definite kernels.  The empirical
arrays are the hollowed Gram matrices
\[
 A_{ij}=k(X_i,X_j)\mathbf 1\{i\ne j\},\qquad
 B_{ij}=\ell(Y_i,Y_j)\mathbf 1\{i\ne j\}.
\]
Then $\Theta_2(k,\ell)$ is the population HSIC, and
\eqref{eq:pair-compression} gives its unbiased U-statistic estimator
\citep{GrettonEtAl2008,SongEtAl2012}.  Bounded kernels automatically satisfy
the standing second-moment condition.  The correspondence between distance
and kernel formulations is discussed by \citet{SejdinovicEtAl2013}.


\section{Regression on pair arrays after U-centering}
\label{sec:edge-regression}

U-centering has two complementary uses.  First, the normalized cross-product
of two residual arrays is an unbiased estimator of a population Hoeffding
pairing, as shown in Section~\ref{sec:pairwise-hoeffding}.  Second, the
same construction gives a regression interpretation.  Treat the entry of one
pair array on each edge as the response and the corresponding entry of a
second array as the predictor.  U-centering both arrays first removes additive
endpoint effects; the slope is then fitted from the residual entries, with one
observation for each unordered pair.

\subsection{Single-predictor regression after residualization}

Let $y,x\in\R^M$ be the response and predictor vectors, let
$W\in\R^{M\times k}$ be a nuisance design matrix of rank $k$, and let $P_W$
be the orthogonal projector onto
$\operatorname{col}(W)^\perp$.  Thus $P_W$ removes from a vector the part
explained by $W$.  Write
\[
 d=\rank(P_W)=M-k.
\]
Provided the adjusted vectors $P_Wy$ and $P_Wx$ are nonzero, their correlation
is
\begin{equation}
 R_W(y,x)=
 \frac{y^\mathsf TP_Wx}
 {\{(y^\mathsf TP_Wy)(x^\mathsf TP_Wx)\}^{1/2}}.
 \label{eq:general-partial-edge-correlation}
\end{equation}

\begin{proposition}[Single-predictor residual regression]
\label{prop:general-edge-regression}
Suppose $d\ge2$, $P_Wx\ne0$, $P_Wy\ne0$, and $|R_W(y,x)|<1$.  In the linear model
\[
 y=W\gamma+\beta x+e,
\]
the ordinary least-squares statistic for testing $H_0:\beta=0$ is
\begin{equation}
 T_W=
 \frac{\sqrt{d-1}\,R_W(y,x)}
 {\{1-R_W(y,x)^2\}^{1/2}},
 \label{eq:general-edge-regression-t}
\end{equation}
and the full-model residual degrees of freedom are $d-1$.
\end{proposition}

This is the Frisch--Waugh--Lovell theorem in correlation form.  First remove
from $y$ and $x$ the parts explained by $W$, leaving $P_Wy$ and $P_Wx$.
Then regress the adjusted response on the adjusted predictor through the
origin and use the usual standard error for a single slope.

\subsection{Dependence as a slope after endpoint residualization}

Let $n\ge4$ and $A,B\in\cV_{n,2}$, with edge vectorizations $a_E$ and
$b_E$, and let $Z$ be the unsigned edge-by-vertex incidence matrix from
Section~\ref{sec:pairwise-anova}.  Consider
\begin{equation}
 b_E=Z\gamma+\beta a_E+e.
 \label{eq:dyadic-regression}
\end{equation}
The term $Z\gamma$ gives each sample label its own additive effect.  Changing
$\gamma_i$ shifts every pair value that contains label $i$.  No separate
intercept is needed: adding the same constant to every endpoint effect shifts
every fitted edge value by a common amount.
By Theorem~\ref{thm:pair-anova}, projecting away these endpoint effects is
exactly the same as U-centering $A$ and $B$.
Suppose that $\widetilde A$ and $\widetilde B$ are nonzero.  Their partial
correlation is then
\begin{equation}
 R_U(A,B)=
 \frac{\inner{\widetilde A}{\widetilde B}_E}
 {\norm{\widetilde A}_E\norm{\widetilde B}_E}.
 \label{eq:u-residual-correlation}
\end{equation}
The residual space after fitting the endpoint effects has dimension
\[
 d_{n,2}=\rank(M_{n,2})=\binom n2-n=\frac{n(n-3)}2,
\]
so, if $|R_U(A,B)|<1$, Proposition~\ref{prop:general-edge-regression} gives
\begin{equation}
 T_{\mathrm{dep}}
 =\frac{\sqrt{d_{n,2}-1}\,R_U(A,B)}
 {\{1-R_U(A,B)^2\}^{1/2}}.
 \label{eq:edge-regression-t}
\end{equation}
We call $T_{\mathrm{dep}}$ the U-centered dependence $t$-statistic.  By
construction, it is the ordinary slope $t$-statistic in
\eqref{eq:dyadic-regression}.  When $A$ and $B$ are Euclidean distance arrays,
$R_U(A,B)$ is the bias-corrected sample distance correlation $R_n^*$, and
$T_{\mathrm{dep}}$ is the corresponding distance-correlation $t$-statistic of
\citet{SzekelyRizzo2013}.
\citet{ZhuZhangYaoShao2020} extend the same studentization to distance- and
reproducing kernel Hilbert space (RKHS)-based dependence measures.  This slope
regression has the residual degrees-of-freedom count
\[
 \binom n2-n-1:
 \qquad
 \text{pair observations}-\text{endpoint effects}-\text{one slope}.
\]

\subsection{The generalized-energy statistic as a slope on the between-group pair indicator}

Let $P$ and $Q$ be probability distributions on a common space $\mathcal Z$.
Let $X,X'$ be independent draws from $P$ and let $Y,Y'$ be independent draws
from $Q$, with all four random variables mutually independent.  For a
symmetric measurable dissimilarity $\rho:\mathcal Z^2\to\R$ for which the
expectations below are finite, define the population generalized-energy
contrast by
\[
 \mathcal E_\rho(P,Q)
 =2\E\rho(X,Y)-\E\rho(X,X')-\E\rho(Y,Y').
\]
The ordinary energy distance is the special case $\rho(x,y)=\norm{x-y}$.

Let $X_1,\ldots,X_n$ and $Y_1,\ldots,Y_m$, with $n,m\ge2$, be independent
samples from $P$ and $Q$, respectively.  Pool the observations as
\[
 U_1=X_1,\ldots,U_n=X_n,
 \qquad
 U_{n+1}=Y_1,\ldots,U_{n+m}=Y_m,
\]
and put $N=n+m$.  Set $G_i=1$ for $i\le n$ and $G_i=0$ for $i>n$ to record
group membership.  For $i<j$, define
\begin{equation}
 C_{ij}=\mathbf 1\{G_i\ne G_j\}=(G_i-G_j)^2,
 \qquad
 A_{ij}=\rho(U_i,U_j),
 \label{eq:between-group-pair-indicator}
\end{equation}
and let $a_E=\operatorname{vec}_E(A)$ and
$c_E=\operatorname{vec}_E(C)$.  Thus $c_E$ indicates which pairs are
between-group pairs.

The following three block averages give an unbiased sample version of the
population contrast:
\[
 \overline A_{XY}=\frac1{nm}\sum_{i\le n<j}A_{ij},
 \quad
 \overline A_{XX}=\frac{2}{n(n-1)}\sum_{i<j\le n}A_{ij},
 \quad
 \overline A_{YY}=\frac{2}{m(m-1)}\sum_{n<i<j\le N}A_{ij},
\]
and, for this fixed choice of $\rho$, the sample generalized-energy statistic is
\[
 \mathcal E_{n,m}=2\overline A_{XY}-\overline A_{XX}-\overline A_{YY}.
\]
This formula uses the dissimilarity convention, under which large cross-group
values contribute positively.  If a positive-definite kernel $k$ is used
directly as a similarity, define $\overline k_{XX}$, $\overline k_{XY}$, and
$\overline k_{YY}$ by replacing $A_{ij}$ with $k(U_i,U_j)$ in the three block
averages above.  The usual unbiased estimator of squared maximum mean
discrepancy (MMD) then has the opposite sign:
\[
 \widehat{\operatorname{MMD}}_{u,k}^2
 =\overline k_{XX}+\overline k_{YY}-2\overline k_{XY}
 =-\{2\overline k_{XY}-\overline k_{XX}-\overline k_{YY}\}.
\]
Thus, if $A_{ij}=k(U_i,U_j)$, the regression slope below is
$-\widehat{\operatorname{MMD}}_{u,k}^2/2$ under the kernel-similarity
convention \citep{SejdinovicEtAl2013}.

To express this contrast as a regression coefficient, define the
block-constant edge vector
$\ell_E=(\ell_{ij}:1\le i<j\le N)$ by
\begin{equation}
 \begin{gathered}
 \ell_{ij}=
 \begin{cases}
 -\{n(n-1)\}^{-1},&1\le i<j\le n,\\[1mm]
 (nm)^{-1},&1\le i\le n<j\le N,\\[1mm]
 -\{m(m-1)\}^{-1},&n<i<j\le N,
 \end{cases}\\[2mm]
 a_{n,m}:=\norm{\ell_E}>0,
 \qquad
 a_{n,m}^2
 =\frac1{nm}+\frac1{2n(n-1)}+\frac1{2m(m-1)}.
 \end{gathered}
 \label{eq:energy-edge-contrast-vector}
\end{equation}
Let $Z_N\in\R^{\binom{N}{2}\times N}$ be the unsigned edge-by-vertex
incidence matrix for the complete graph on the $N$ pooled labels, and let
$M_{N,2}$ be its residual projector,
defined as in \eqref{eq:pair-projection}.  The common-endpoint edge regression is
\begin{equation}
 a_E=Z_N\gamma+\theta c_E+e.
 \label{eq:simple-pooled-group-regression}
\end{equation}

\begin{proposition}[The generalized-energy statistic is a slope on the between-group pair indicator]
\label{prop:energy-group-slope}
The residualized between-group pair indicator satisfies
\begin{equation}
 M_{N,2}c_E=\frac{\ell_E}{a_{n,m}^2}.
 \label{eq:ucentered-group-indicator}
\end{equation}
Consequently, the least-squares slope in
\eqref{eq:simple-pooled-group-regression} is
\begin{equation}
 \widehat\theta
 =\ell_E^\mathsf Ta_E
 =\overline A_{XY}-\frac{\overline A_{XX}+\overline A_{YY}}2
 =\frac{\mathcal E_{n,m}}2.
 \label{eq:energy-as-group-slope}
\end{equation}
Equivalently, if $\widetilde A$ and $\widetilde C$ denote global U-centering
on all $N$ labels, then
\begin{equation}
 \mathcal E_{n,m}
 =2a_{n,m}^2\inner{\widetilde A}{\widetilde C}_E.
 \label{eq:energy-as-ucentered-label-covariance}
\end{equation}
\end{proposition}

The fitted slope compares average between-group and within-group
dissimilarities after additive endpoint effects have been removed.  The
explicit projection of $c_E$ is proved in
Appendix~\ref{app:two-sample-nesting}.

\subsection{Restricted and fully interacted studentizations}

We now compare two studentizations, meaning two ways of dividing the fitted
slope by an estimated standard error.
In the common-endpoint model \eqref{eq:simple-pooled-group-regression}, each
pooled observation has the same endpoint effect in its within- and
between-group edges.  After fitting the generalized-energy slope
$\widehat\theta$, the residual sum of
squares and its degrees of freedom are
\begin{equation}
 \mathrm{RSS}_R
 =a_E^\mathsf TM_{N,2}a_E-\frac{\widehat\theta^2}{a_{n,m}^2},
 \qquad
 \nu_R=\binom N2-N-1.
 \label{eq:restricted-two-sample-rss}
\end{equation}
Provided $\mathrm{RSS}_R>0$, the slope statistic for this common-endpoint
model is
\begin{equation}
 T_R=\frac{\widehat\theta}
 {a_{n,m}\sqrt{\mathrm{RSS}_R/\nu_R}}.
 \label{eq:restricted-energy-t}
\end{equation}
We call the regression the restricted model and the resulting standardization
of $\widehat\theta$ the restricted studentization.  A closely related pooled
studentization is due to \citet{GaoShao2023};
Remark~\ref{rem:gao-shao} gives the exact relation.

For the remainder of this subsection, assume $n,m\ge4$.  We now introduce a
fully interacted regression whose slope statistic
reproduces the pooled studentization of \citet{ChakrabortyZhang2021}.  This
regression uses a more flexible model for endpoint effects.  Put
$Q_{ij}=G_i+G_j-1$, which equals $1$, $0$, and $-1$ on $XX$, $XY$, and $YY$
edges, respectively, and fit
\begin{equation}
 \begin{split}
 A_{ij}={}&\eta+\delta Q_{ij}+\theta C_{ij}\\
 &+(1-C_{ij})(\alpha_i+\alpha_j)
   +C_{ij}(\beta_i+\beta_j)+e_{ij},
 \qquad i<j.
 \end{split}
 \label{eq:stacked-two-sample-edge-model}
\end{equation}
The three block intercepts are $\eta+\delta$ for $XX$, $\eta+\theta$ for
$XY$, and $\eta-\delta$ for $YY$.  Thus $\delta$ distinguishes the two
within-group blocks, while $\theta$ compares the between-group block with
their average.
For identifiability, the $\alpha$ effects sum to zero separately in the two
groups, and so do the $\beta$ effects.  The $\alpha_i$ are endpoint effects
for within-group edges, whereas the $\beta_i$ are endpoint effects for
between-group edges.  Thus this model allows the effect of an observation to
depend on the type of edge in which it appears.

Partition $A$ into two symmetric hollow within-group arrays and one
rectangular between-group array:
\[
 A^{XX}=(A_{ij})_{1\le i,j\le n},\qquad
 A^{XY}=(A_{i,n+j})_{\substack{1\le i\le n\\1\le j\le m}},\qquad
 A^{YY}=(A_{n+i,n+j})_{1\le i,j\le m}.
\]
Also define
\[
 \Pi_s=I_s-\frac1s\one_s\one_s^\mathsf T.
\]
Write $\norm{\cdot}_F$ for the Frobenius norm.

\begin{proposition}[Nested edge regressions for two-sample homogeneity]
\label{prop:energy-anova-t}
Let $n,m\ge4$.  In model \eqref{eq:stacked-two-sample-edge-model}:
\begin{enumerate}
\item the slope is again $\widehat\theta=\mathcal E_{n,m}/2$;
\item the residual sum of squares is
\begin{equation}
 \mathrm{RSS}_F=\mathrm{RSS}_{XY}+\mathrm{RSS}_{XX}+\mathrm{RSS}_{YY},
 \label{eq:stacked-full-rss}
\end{equation}
where
\begin{align*}
 \mathrm{RSS}_{XY}
 &=\norm{\Pi_nA^{XY}\Pi_m}_F^2,\\
 \mathrm{RSS}_{XX}
 &=\operatorname{vec}_E(A^{XX})^\mathsf T
   M_{n,2}\operatorname{vec}_E(A^{XX}),\\
 \mathrm{RSS}_{YY}
 &=\operatorname{vec}_E(A^{YY})^\mathsf T
   M_{m,2}\operatorname{vec}_E(A^{YY});
\end{align*}
\item the residual degrees of freedom are
\begin{equation}
 \nu_F
 =(n-1)(m-1)+\frac{n(n-3)}2+\frac{m(m-3)}2
 =\binom N2-2N+1;
 \label{eq:stacked-two-sample-df}
\end{equation}
\item let $S_{n,m}$ denote the pooled variance estimator of
\citet{ChakrabortyZhang2021}, whose blockwise formula is reproduced in
\eqref{eq:cz-pooled-variance}.  It satisfies
$S_{n,m}=4\mathrm{MS}_F$, where
$\mathrm{MS}_F=\mathrm{RSS}_F/\nu_F$.  Consequently, provided
$\mathrm{RSS}_F>0$, the ordinary slope statistic is
\begin{equation}
 T_F=\frac{\widehat\theta}{a_{n,m}\sqrt{\mathrm{MS}_F}}
 =\frac{\mathcal E_{n,m}}{a_{n,m}\sqrt{S_{n,m}}},
 \label{eq:energy-anova-t}
\end{equation}
which, under the same dissimilarity convention, is
exactly their studentized statistic;
\item the restricted model is nested in the fully interacted model and
\begin{equation}
 \mathrm{RSS}_R=\mathrm{RSS}_F+\mathrm{SS}_{\mathrm{int}},
 \qquad
 \mathrm{SS}_{\mathrm{int}}\ge0,
 \qquad
 \nu_R-\nu_F=N-2,
 \label{eq:nested-rss-decomposition}
\end{equation}
where $\mathrm{SS}_{\mathrm{int}}$ is the extra sum of squares explained by
the endpoint-by-pair-type interactions.
\end{enumerate}
\end{proposition}

Thus the two procedures differ only in the nuisance fit used to form the
residual mean square.

\begin{remark}[Connection with Gao and Shao]
\label{rem:gao-shao}
Set $d=d_{N,2}$ and $Q_0=a_E^\mathsf TM_{N,2}a_E$.  Under the
constant-diagonal condition of \citet{GaoShao2023}, and after aligning the
dissimilarity convention, let $\widetilde A$ be the globally U-centered
hollow pooled array and write $a_0^k$ for the common diagonal value.  Their
centered off-diagonal entries equal
$\widetilde A_{st}-a_0^k/(N-1)$, and their explicit diagonal correction
cancels the squared common shift.  Their pooled estimator therefore reduces
to
\[
 \frac{1}{N(N-3)}\sum_{s\ne t}\widetilde A_{st}^{\,2}
 =\frac{Q_0}{d},
\]
because the ordered-pair sum equals $2Q_0$ and $N(N-3)=2d$.
If $Q_0>0$, their studentized statistic can therefore be written as
\[
 T_{\mathrm{GS}}
 =\frac{\widehat\theta}{a_{n,m}\sqrt{Q_0/d}}.
\]
Because $\mathrm{RSS}_R=Q_0-\widehat\theta^2/a_{n,m}^2$ and
$\nu_R=d-1$, whenever $\mathrm{RSS}_R>0$,
\[
 T_R=T_{\mathrm{GS}}
 \left\{\frac{d-1}{d-T_{\mathrm{GS}}^2}\right\}^{1/2}.
\]
Thus the two statistics are signed monotone transformations and give the same
two-sided ordering.  Gao and Shao use the residual mean square before fitting
the between-group slope, whereas $T_R$ uses the ordinary residual mean square
after the slope has been fitted.
\end{remark}

\begin{remark}[Fixed-sample, high-dimensional limits]
When the sample sizes are fixed and the ambient dimension grows,
\citet{ZhuZhangYaoShao2020} and \citet{ChakrabortyZhang2021} show, under their
respective null hypotheses and assumptions, that $T_{\mathrm{dep}}$ and $T_F$
converge to Student $t$ distributions.
\end{remark}

\section{Higher-order U-centering}
\label{sec:higher-order-u-centering}

We now extend pairwise residualization to arrays indexed by $r$-subsets.  The
endpoint effects are replaced by all lower-order subset effects, and the zero
row sums become zero $(r-1)$-way margins.

\subsection{\texorpdfstring{{The hierarchy for triple-indexed arrays}}{The hierarchy for triple-indexed arrays}}

Assume $n\ge6$, and consider a symmetric array $A_{ijk}$ indexed by unordered
triples of distinct labels.  Its lower-order hierarchical additive fit is
\[
 A_{ijk}=\mu+\alpha_i+\alpha_j+\alpha_k
 +\beta_{ij}+\beta_{ik}+\beta_{jk}+\varepsilon_{ijk}.
\]
The parametrization is redundant: individual coefficients are not unique,
but the fitted array and its least-squares residual are.  Suitable linear
constraints may be imposed to select a unique parametrization without
changing either quantity.

Here $\varepsilon_{ijk}$ denotes the deterministic least-squares residual,
not a stochastic error term.  Its normal equations are the zero pair-margin
conditions
\[
 \sum_{k\notin\{i,j\}}\varepsilon_{ijk}=0
 \qquad\text{for every }\{i,j\}\in\binom{[n]}2.
\]
Thus U-centering for triple-indexed arrays removes every effect involving fewer
than three labels.  The condition $n\ge6$ is the $r=3$ specialization of the
general condition $n\ge2r$.

\subsection{Arrays, margins, and the lower-order ANOVA space}

Fix $r\ge1$ and $n\ge2r$, and let
\[
 \cI_{n,r}=\binom{[n]}r,
 \qquad
 \cV_{n,r}=\R^{\cI_{n,r}}.
\]
An element $A\in\cV_{n,r}$ is an $r$-subset array evaluated only at
distinct sample labels.  For $A,B\in\cV_{n,r}$, define the inner product by
  \begin{equation*}
 \inner{A}{B}=\sum_{I\in\cI_{n,r}}A_I B_I.
\end{equation*}
For $0\le j\le r$, with $\binom{[n]}0=\{\varnothing\}$, define the margin map
$D_{j,r}:\cV_{n,r}\to\R^{\binom{[n]}{j}}$ by
  \begin{equation*}
 (D_{j,r}A)(J)
 =\sum_{\substack{I\in\cI_{n,r}\\I\supseteq J}}A_I,
 \qquad J\in\binom{[n]}j.
\end{equation*}
For $j=0$, this is the grand sum $D_{0,r}A=\sum_I A_I$.  If
$f:\binom{[n]}j\to\R$, the adjoint map is
  \begin{equation*}
 (D_{j,r}^\ast f)(I)
 =\sum_{\substack{J\subseteq I\\|J|=j}}f(J).
\end{equation*}
Thus $D_{j,r}$ collects $j$-set margins, whereas $D_{j,r}^\ast$ spreads a
function on $j$-sets back to the $r$-sets containing them.

The full lower-order ANOVA fit is
\begin{equation}\label{eq:hyper-anova}
 A_I
 =\sum_{j=0}^{r-1}\ \sum_{\substack{J\subseteq I\\|J|=j}}
   \alpha_J+\varepsilon_I.
\end{equation}
Although this parametrization is redundant, its fitted-value space is
  \begin{equation*}
 \cL_{n,r}=\range(D_{r-1,r}^\ast).
\end{equation*}
To see this, let $f$ be a function on $j$-sets with $j<r-1$.  Then
  \begin{equation*}
 D_{r-1,r}^\ast D_{j,r-1}^\ast f
 =(r-j)D_{j,r}^\ast f.
\end{equation*}
For a fixed $r$-set $I$, each $j$-subset $J\subseteq I$ is contained in
exactly $r-j$ of the $(r-1)$-subsets of $I$, which proves the identity.
Therefore, all lower-order fitted effects are contained in
$\range(D_{r-1,r}^\ast)$.  Conversely, the $j=r-1$ term in
\eqref{eq:hyper-anova} spans this range.  Hence the full fitted-value space is
exactly $\cL_{n,r}=\range(D_{r-1,r}^\ast)$.  This nesting property of inclusion
matrices is also discussed in \citet{Gottlieb1966}.  The residual, or pure
$r$-way interaction, space is
  \[
 \cH_{n,r}=\cL_{n,r}^{\perp}=\ker(D_{r-1,r}).
\]
Indeed, since $\cL_{n,r}=\range(D_{r-1,r}^{\ast})$, for any
$H\in\cV_{n,r}$,
\[
\begin{aligned}
 H\in\cL_{n,r}^{\perp}
 &\quad\Longleftrightarrow\quad
 \inner{H}{D_{r-1,r}^{\ast}f}=0
 \quad\text{for every }f\in\R^{\binom{[n]}{r-1}}\\
 &\quad\Longleftrightarrow\quad
 \inner{D_{r-1,r}H}{f}=0
 \quad\text{for every }f
 \quad\Longleftrightarrow\quad
 D_{r-1,r}H=0.
\end{aligned}
\]
Hence $H\in\cH_{n,r}$ if and only if
  \[
 \sum_{i\notin J}H_{J\cup\{i\}}=0
 \qquad\text{for every }J\in\binom{[n]}{r-1}.
\]
\subsection{An explicit higher-order projection}
We seek coefficients $c_{n,r,j}$ for a linear combination of the margin
operators $D_{j,r}^\ast D_{j,r}$.  Since $D_{r,r}$ is the identity, set
$c_{n,r,r}=1$ to retain the raw entry, and choose the remaining coefficients
so that every $(r-1)$-margin vanishes.
The margin count carried out in
Appendix~\ref{app:higher-projection-proof} yields the recurrence
 \[
 c_{n,r,j+1}
 =-\frac{n-r-j+1}{r-j}c_{n,r,j},
 \qquad 0\le j<r.
\]
Solving this recurrence backward from $c_{n,r,r}=1$ gives
  \[
 c_{n,r,j}
 =\frac{(-1)^{r-j}}{\binom{n-r-j+1}{r-j}}.
\]
\begin{definition}[Higher-order U-centering]
\label{def:higher-u}
For $A\in\cV_{n,r}$, define
\begin{equation}\label{eq:higher-u-operator}
 \cU_{n,r}A
 =\sum_{j=0}^r c_{n,r,j}D_{j,r}^\ast D_{j,r}A.
\end{equation}
Equivalently, for $I\in\cI_{n,r}$,
\begin{equation}\label{eq:higher-u-entry}
 (\cU_{n,r}A)_I
 =\sum_{j=0}^r
 \frac{(-1)^{r-j}}{\binom{n-r-j+1}{r-j}}
 \sum_{\substack{J\subseteq I\\|J|=j}}
 \sum_{\substack{I'\in\cI_{n,r}\\I'\supseteq J}}A_{I'}.
\end{equation}
\end{definition}

For $r=1$, \eqref{eq:higher-u-entry} is ordinary sample centering.  For
$r=2$ and $n\ge4$, it becomes
\[
 (\cU_{n,2}A)_{ij}
 =A_{ij}-\frac{A_{i\cdot}+A_{j\cdot}}{n-2}
       +\frac{A_{\cdot\cdot}}{(n-1)(n-2)},
\]
which is pairwise U-centering.

\begin{samepage}
\begin{theorem}[Higher-order U-centering is the ANOVA projection]
\label{thm:higher-projection}
Let $n\ge2r$.  The operator $\cU_{n,r}$ is the orthogonal projection of
$\cV_{n,r}$ onto the residual space
$\cH_{n,r}=\ker(D_{r-1,r})$.  Equivalently, for every
$A\in\cV_{n,r}$, $\cU_{n,r}A$ is the least-squares residual obtained after
fitting every ANOVA effect involving fewer than $r$ labels.  In particular,
\begin{enumerate}
\item $D_{r-1,r}(\cU_{n,r}A)=0$ for every $A$;
\item $\cU_{n,r}A=A$ for every $A\in\cH_{n,r}$;
\item $\cU_{n,r}$ is self-adjoint and idempotent;
\item $\ker(\cU_{n,r})=\cL_{n,r}$; and
\item the residual degrees of freedom are given by
\begin{equation}\label{eq:general-df}
 d_{n,r}:=\dim(\cH_{n,r})
 =\binom nr-\binom n{r-1}
 =\binom nr\frac{n-2r+1}{n-r+1}.
\end{equation}
\end{enumerate}
\end{theorem}
\end{samepage}

\noindent The proof is given in
Appendix~\ref{app:higher-projection-proof}.

\subsection{\texorpdfstring{{The explicit projection for triple-indexed arrays}}{The explicit projection for triple-indexed arrays}}

Assume $n\ge6$.  Let $A_{ijk}$ be symmetric and defined only for distinct
$i,j,k$.  The margins for a triple-indexed array are
\begin{align*}
 M_2(i,j)&=\sum_{k\notin\{i,j\}}A_{ijk},\\
 M_1(i)&=\sum_{\substack{j<k\\j,k\ne i}}A_{ijk},\\
 M_0&=\sum_{i<j<k}A_{ijk}.
\end{align*}
Then, for distinct $i,j,k$,
\begin{equation}\label{eq:r3-formula}
\begin{split}
 (\cU_{n,3}A)_{ijk}
 ={}&A_{ijk}
 -\frac{M_2(i,j)+M_2(i,k)+M_2(j,k)}{n-4}\\
 &+\frac{M_1(i)+M_1(j)+M_1(k)}{\binom{n-3}{2}}
 -\frac{M_0}{\binom{n-2}{3}}.
\end{split}
\end{equation}
The alternating terms follow the ANOVA hierarchy: retain the raw triple,
subtract the appropriately scaled pair margins, add back the scaled one-label
margins, and finally adjust for the grand margin.
The resulting pair-margin equations are the normal equations for the
hierarchical residual.  In particular, every pairwise slice has zero sum:
  \begin{equation*}
 \sum_{k\notin\{i,j\}}(\cU_{n,3}A)_{ijk}=0.
\end{equation*}
The residual degrees of freedom are
  \begin{equation*}
 d_{n,3}=\binom n3-\binom n2
 =\frac{n(n-1)(n-5)}{6}.
\end{equation*}

\subsection{Higher-order subset regression}
\label{sec:higher-order-subset-regression}

The pairwise construction extends by replacing edges with $r$-subsets and
endpoint effects with all lower-order subset effects.  Fix an ordering of
$\cI_{n,r}$ and identify
$A,B\in\cV_{n,r}$ with the coordinate vectors
$\mathbf a=(A_I:I\in\cI_{n,r})$ and
$\mathbf b=(B_I:I\in\cI_{n,r})$ in $\R^{\binom nr}$.  By
\eqref{eq:general-df}, $\dim(\cL_{n,r})=\binom n{r-1}$.  Choose a
full-column-rank matrix
\[
 W_{<r}\in\R^{\binom nr\times\binom n{r-1}}
\]
whose column space is the lower-order ANOVA space $\cL_{n,r}$.  For
$\gamma\in\R^{\binom n{r-1}}$, consider
\begin{equation}
 \mathbf b=W_{<r}\gamma+\beta\mathbf a+\mathbf e,
 \label{eq:higher-order-subset-regression}
\end{equation}
where $\mathbf e$ is the residual vector.  Under this coordinate
identification, $P_{W_{<r}}=\cU_{n,r}$ because both
are the orthogonal projector onto $\cL_{n,r}^{\perp}$.  Thus the
regression first U-centers both arrays, removing every effect involving fewer
than $r$ labels, and then fits the slope between the adjusted arrays.

\begin{corollary}[Single-predictor regression on $r$-subsets]
\label{cor:higher-order-subset-regression}
Let $n\ge2r$ and $d_{n,r}\ge2$.  Suppose $\cU_{n,r}A$ and
$\cU_{n,r}B$ are nonzero, and define
\begin{equation}
 R_{n,r}^{\mathrm{sub}}(A,B)
 =\frac{\inner{\cU_{n,r}A}{\cU_{n,r}B}}
 {\norm{\cU_{n,r}A}\norm{\cU_{n,r}B}}.
 \label{eq:higher-order-residual-correlation}
\end{equation}
If $|R_{n,r}^{\mathrm{sub}}(A,B)|<1$, the ordinary least-squares statistic
for testing $H_0:\beta=0$ in
\eqref{eq:higher-order-subset-regression} is
\begin{equation}
 T_{n,r}^{\mathrm{slope}}
 =\frac{\sqrt{d_{n,r}-1}\,R_{n,r}^{\mathrm{sub}}(A,B)}
 {\{1-(R_{n,r}^{\mathrm{sub}}(A,B))^2\}^{1/2}},
 \qquad
 d_{n,r}=\binom nr-\binom n{r-1},
 \label{eq:higher-order-subset-t}
\end{equation}
and its full-model residual degrees of freedom are $d_{n,r}-1$.
\end{corollary}

The corollary follows from Proposition~\ref{prop:general-edge-regression} with
$M=\binom nr$, $P_W=\cU_{n,r}$, and $d=d_{n,r}$, together with
Theorem~\ref{thm:higher-projection}.

The same construction extends to several subset-level predictors: U-center
the response and every predictor column, then fit the resulting ordinary
multiple regression, provided the residualized predictor columns are linearly
independent.

\section{Population pairings and projection representation}
\label{sec:population-pairings}

\subsection{General Hoeffding components}\label{sec:general_hoeffding}

We now introduce the population counterpart of the sample-label ANOVA\@.
Fix $r\ge1$, and write $P_X$ for the law of $X$.  Let
$O_i=(X_i,Y_i)$, $i\ge1$, be i.i.d.\ copies of $O=(X,Y)$.  For
$I=\{i_1<\cdots<i_r\}$, write $X_I=(X_{i_1},\ldots,X_{i_r})$ and define
$Y_I$ analogously.  Let $a:\mathcal X^r\to\R$ and
$b:\mathcal Y^r\to\R$ be symmetric kernels satisfying
\[
 \E\{a(X_1,\ldots,X_r)^2\}<\infty,
 \qquad
 \E\{b(Y_1,\ldots,Y_r)^2\}<\infty.
\]
Unlike the sample-label ANOVA of Section~\ref{sec:higher-order-u-centering}, this decomposition takes place in
a population $L^2$ space and is defined by conditional expectation.  We use
the usual Hoeffding components.  To define them, write
\[
 g_j(x_1,\ldots,x_j)
 =\E\{a(x_1,\ldots,x_j,X_{j+1},\ldots,X_r)\},
 \qquad 0\le j\le r,
\]
where $g_0=\E a(X_1,\ldots,X_r)$ and $g_r=a$.  For
$S=\{s_1<\cdots<s_\ell\}\subseteq[j]$, write
$x_S=(x_{s_1},\ldots,x_{s_\ell})$.  Set $a_0=g_0$, interpret
$a_0(x_\varnothing)$ as $a_0$, and define recursively
\[
 a_j(x_1,\ldots,x_j)
 =g_j(x_1,\ldots,x_j)
 -\sum_{\ell=0}^{j-1}
   \sum_{\substack{S\subseteq[j]\\|S|=\ell}}a_\ell(x_S).
\]
Equivalently,
\[
 g_j(x_1,\ldots,x_j)
 =\sum_{S\subseteq[j]}a_{|S|}(x_S).
\]
Taking $j=r$ gives the Hoeffding decomposition
\begin{equation}\label{eq:general-hoeffding-decomp}
 a(X_I)=\sum_{S\subseteq I}a_{|S|}(X_S),
 \qquad I\in\binom{[n]}r.
\end{equation}
A symmetric $j$-argument kernel is called \emph{canonical} (or completely
degenerate) if integrating out any one argument gives zero for
$P_X^{j-1}$-almost every choice of the remaining arguments.  The components
in the standard Hoeffding decomposition above are symmetric and canonical;
in particular,
\begin{equation}\label{eq:general-canonical}
 \E\{a_j(x_1,\ldots,x_{j-1},X_j)\}=0,
 \qquad j\ge1.
\end{equation}
See \citet{Hoeffding1948} and \citet{Serfling1980} for the standard
derivation and canonicality properties.
We define $b_0,\ldots,b_r$ analogously.

For $0\le j\le r$, we define the bilinear interaction pairings
\begin{equation*}
 \Theta_j(a,b)
 =\E\{a_j(X_1,\ldots,X_j)b_j(Y_1,\ldots,Y_j)\},
\end{equation*}
where $\Theta_0=a_0b_0$.  The target of interest is the $r$th
Hoeffding-component pairing $\Theta_r(a,b)$.
If $X\perp Y$, then for every $j\ge1$ the two canonical components are
independent and have mean zero, so $\Theta_j(a,b)=0$.

\subsection{Population moments indexed by overlap size}

For $0\le s\le r$, take $2r-s$ i.i.d.\ paired observations and set
\[
 I=\{1,\ldots,r\},
 \qquad
 J_s=\{1,\ldots,s,r+1,\ldots,2r-s\}.
\]
Both $I$ and $J_s$ contain $r$ labels, and $|I\cap J_s|=s$.  We define
\[
 \eta_s(a,b)=\E\{a(X_I)b(Y_{J_s})\}.
\]
Equivalently,
\[
 \eta_s(a,b)
 =\E\left[
 a(X_1,\ldots,X_r)
 b(Y_1,\ldots,Y_s,Y_{r+1},\ldots,Y_{2r-s})
 \right].
\]
The subscripts denote sample labels rather than argument positions: the two
kernel evaluations share the first $s$ paired observations and use distinct
observations for their remaining arguments.  The following triangular identity
gives a bilinear form of the classical
covariance decomposition for Hoeffding components
\citep{Hoeffding1948,Serfling1980}.
When $a=b=h$, it reduces to the familiar relation between an
overlap covariance and the Hoeffding variance components.  In the present
notation, it also provides the link to the finite-sample projection below.


\begin{theorem}
\label{thm:overlap-anova}
For every $0\le s\le r$,
\begin{equation}\label{eq:eta-theta-triangular}
 \eta_s(a,b)=\sum_{j=0}^s\binom sj\Theta_j(a,b).
\end{equation}
Consequently,
\begin{equation}\label{eq:theta-binomial-inversion}
 \Theta_r(a,b)
 =\sum_{s=0}^r(-1)^{r-s}\binom rs\eta_s(a,b).
\end{equation}
\end{theorem}

\begin{proof}
Expand the two raw kernels using
\eqref{eq:general-hoeffding-decomp}.  Consider a term
$a_{|S|}(X_S)b_{|T|}(Y_T)$ with $S\subseteq I$ and $T\subseteq J_s$.
If $S\ne T$, some sample label belongs to exactly one of the two sets.  On
conditioning on all other labels, canonicality
\eqref{eq:general-canonical} makes the expectation zero.  A nonzero term must
therefore have $S=T\subseteq I\cap J_s=[s]$.  For each $j$, there are
$\binom sj$ shared subsets of size $j$, and each contributes
$\Theta_j(a,b)$.  This proves \eqref{eq:eta-theta-triangular};
\eqref{eq:theta-binomial-inversion} is ordinary binomial inversion.
\end{proof}
For $r=2$, Theorem~\ref{thm:overlap-anova} reduces to
Proposition~\ref{prop:pair-overlap}.
More generally, the overlap-$s$ moment contains Hoeffding components through
order $s$, and binomial inversion isolates the order-$r$ component.


\subsection{\texorpdfstring{{A spectral example based on four-argument matching kernels}}{A spectral example based on four-argument matching kernels}}\label{sec:example}
Let $k$ and $\ell$ be bounded measurable positive-definite kernels with
separable RKHSs $\mathcal F$ and $\mathcal G$, respectively.  For four
arguments, write $k_{ij}=k(x_i,x_j)$ and
$\ell_{ij}=\ell(y_i,y_j)$, and define the observable symmetric matching
kernels
\begin{align}
 m_k(x_1,x_2,x_3,x_4)
 &=\frac13\{k_{12}k_{34}+k_{13}k_{24}+k_{14}k_{23}\},\label{eq:matching-k}\\
 m_\ell(y_1,y_2,y_3,y_4)
 &=\frac13\{\ell_{12}\ell_{34}+\ell_{13}\ell_{24}
                  +\ell_{14}\ell_{23}\}.\label{eq:matching-l}
\end{align}
The three terms correspond to the three perfect matchings of four labels. Let $k_c$ and $\ell_c$ denote the population-centered kernels.  For example,
\[
 k_c(x,x')
 =k(x,x')-\E k(x,X)-\E k(X,x')+\E k(X,X'),
\]
where $X$ and $X'$ are independent with distribution $P_X$; define
$\ell_c$ analogously.  Write
$k_{c,ij}=k_c(x_i,x_j)$ and $\ell_{c,ij}=\ell_c(y_i,y_j)$.  The
fourth Hoeffding components of the raw matching kernels are
\begin{align*}
 (m_k)_4(x_1,x_2,x_3,x_4)
 &=\frac13\{k_{c,12}k_{c,34}+k_{c,13}k_{c,24}
                  +k_{c,14}k_{c,23}\},\\
 (m_\ell)_4(y_1,y_2,y_3,y_4)
 &=\frac13\{\ell_{c,12}\ell_{c,34}+\ell_{c,13}\ell_{c,24}
                  +\ell_{c,14}\ell_{c,23}\}.
\end{align*}
Appendix~\ref{app:rkhs} verifies these projections.
Let
\[
 \phi_c(x)=k(x,\cdot)-\E k(X,\cdot),
 \qquad
 \psi_c(y)=\ell(y,\cdot)-\E \ell(Y,\cdot)
\]
be the centered feature maps.  For $u\in\mathcal F$ and $v,g\in\mathcal G$,
define the rank-one operator $u\otimes v:\mathcal G\to\mathcal F$ by
$(u\otimes v)g=u\inner{v}{g}_{\mathcal G}$.  Define the cross-covariance
operator
\[
 C_{XY}=\E\{\phi_c(X)\otimes\psi_c(Y)\}:\mathcal G\to\mathcal F.
\]
Let $C_{XY}^{\ast}$ denote its adjoint, characterized by
\[
 \inner{C_{XY}g}{f}_{\mathcal F}
 =\inner{g}{C_{XY}^{\ast}f}_{\mathcal G}.
\]
For any orthonormal basis $(f_q)$ of $\mathcal G$, define the
Hilbert--Schmidt norm by
\[
 \norm{C_{XY}}_{\mathrm{HS}}^2
 =\sum_q\norm{C_{XY}f_q}_{\mathcal F}^2,
\]
and let $s_1,s_2,\ldots$ be the singular values of $C_{XY}$.  Then population
HSIC is
\begin{equation}\label{eq:H2-spectrum}
 H_2:=\Theta_2(k,\ell)
 =\norm{C_{XY}}_{\mathrm{HS}}^2
 =\sum_{j\ge1}s_j^2;
\end{equation}
see \citet{GrettonEtAl2005}.  The pairing of the fourth Hoeffding
components of the matching kernels is
\[
 H_4:=\Theta_4(m_k,m_\ell)
 =\E\{(m_k)_4(X_1,\ldots,X_4)
       (m_\ell)_4(Y_1,\ldots,Y_4)\}.
\]

Expanding the two three-term averages gives nine pairs of matchings.  For each
of the three coincident matching pairs, the two matched edges involve
disjoint groups of observations, so the expectation factors as $H_2^2$.  For
each of the six different matching pairs, their union is a four-cycle.  By
relabeling, all six four-cycle terms have the same expectation, and
Appendix~\ref{app:rkhs} shows that
\[
 \E\{k_{c,12}k_{c,34}\ell_{c,13}\ell_{c,24}\}
 =\tr\{(C_{XY}^{\ast}C_{XY})^2\}
 =\sum_{j\ge1}s_j^4.
\]
Combining the three coincident and six four-cycle terms with their common
factor $1/9$ gives
\begin{equation}\label{eq:H4-spectrum}
 \boxed{
 H_4=\frac13H_2^2+\frac23\sum_{j\ge1}s_j^4.
 }
\end{equation}
Both terms on the right are nonnegative, so $H_4=0$ if and only if
$H_2=0$, or equivalently $C_{XY}=0$.  Under the kernel conditions of
\citet{SzaboSriperumbudur2018}, this is further equivalent to $X\perp Y$.
Equation~\eqref{eq:H4-spectrum} recovers the fourth spectral moment from the
two population pairings:
\[
 S_4:=\sum_{j\ge1}s_j^4=\frac{3H_4-H_2^2}{2}.
\]
When $H_2>0$, the corresponding spectral participation ratio is
\begin{equation}\label{eq:spectral-participation}
 d_{\mathrm{eff}}
 :=\frac{(\sum_j s_j^2)^2}{\sum_j s_j^4}
 =\frac{2H_2^2}{3H_4-H_2^2}.
\end{equation}
This ratio is invariant to the choice of bases in the two RKHSs, but it
depends on the chosen kernels.  It should therefore be interpreted as the
effective number of dependence directions in the chosen RKHS geometry, not
as an intrinsic dimension of the joint distribution.  At independence,
$H_2=H_4=0$ and the ratio is undefined.

The distinction can be seen at fixed $H_2=\omega^2>0$.  If the spectrum is
rank one, with $s_1=\omega$, then
\[
 H_4=\omega^4,\qquad d_{\mathrm{eff}}=1.
\]
If instead $s_1=\cdots=s_q=\omega/\sqrt q$ and all remaining singular values
vanish, ordinary HSIC is still $\omega^2$, whereas
\[
 H_4=\omega^4\left(\frac13+\frac{2}{3q}\right),
 \qquad d_{\mathrm{eff}}=q.
\]
Thus $H_2$ records total squared cross-covariance, while $H_4$ distinguishes
whether that total is concentrated in one spectral direction or spread
across several directions.



\subsection{The direct overlap estimator}

We now turn from population interpretation to finite-sample estimation.
Assume $n\ge2r$.  For $I,J\in\cI_{n,r}$, let $A_I=a(X_I)$ and
$B_J=b(Y_J)$.  For $0\le s\le r$, the number of ordered pairs $(I,J)$ with
$|I\cap J|=s$ is
\begin{equation}\label{eq:N-s}
 N_{n,r,s}
 =\binom nr\binom rs\binom{n-r}{r-s}.
\end{equation}
Here $I$ and $J$ are themselves unordered subsets, but $(I,J)$ is an ordered
pair: choose $I$, choose the $s$ labels it shares with $J$, and then choose
the remaining $r-s$ labels of $J$.
Define
\begin{equation}\label{eq:eta-s-hat}
 \widehat\eta_{n,r,s}
 =\frac{1}{N_{n,r,s}}
  \sum_{\substack{I,J\in\cI_{n,r}\\|I\cap J|=s}}A_I B_J.
\end{equation}
By exchangeability, $\E\widehat\eta_{n,r,s}=\eta_s(a,b)$.  We therefore define
the direct unbiased estimator by
\begin{equation}\label{eq:direct-general}
 T_{n,r}^{\mathrm{dir}}
 =\sum_{s=0}^r(-1)^{r-s}\binom rs
  \widehat\eta_{n,r,s}.
\end{equation}
The term $s=0$ alone contains on the order of $n^{2r}$ products, so a literal
evaluation costs $O(n^{2r})$ arithmetic operations.  This count treats the
arrays $A$ and $B$ as already formed; evaluating the kernels on the
$\binom nr$ subsets is a separate, kernel-dependent cost.

Each $\widehat\eta_{n,r,s}$ is a U-statistic of order $2r-s$ after
symmetrization.  A U-statistic of order $q<2r$ can always be lifted to order
$2r$ by averaging its kernel over the $q$-subsets of $2r$ arguments.
Therefore \eqref{eq:direct-general} admits an order-$2r$ U-statistic
representation.  For $r=2$, this is the familiar fourth-order representation;
for $r=3$, it is a sixth-order representation.


\subsection{Direct evaluation by subset-margin inversion}
We first derive the $O(n^r)$ scalar evaluation directly, without using the
projection identity.  Define the unnormalized overlap sums
\begin{equation*}
 \Sigma_s(A,B)
 =\sum_{\substack{I,J\in\cI_{n,r}\\|I\cap J|=s}}A_I B_J,
 \qquad 0\le s\le r,
\end{equation*}
and, for $0\le k\le r$, the cross-products of the $k$-way margins
\begin{equation}\label{eq:Ck-margin}
 C_k(A,B)
 =\sum_{K\in\binom{[n]}k}
   (D_{k,r}A)(K)(D_{k,r}B)(K).
\end{equation}
For $k=0$, recall that $(D_{0,r}A)(\varnothing)=\sum_I A_I$.

\begin{proposition}
\label{prop:margin-inversion}
For $0\le k\le r$,
\begin{equation}\label{eq:Ck-Sigma}
 C_k(A,B)=\sum_{s=k}^r\binom{s}{k}\Sigma_s(A,B).
\end{equation}
Consequently,
\begin{equation}\label{eq:Sigma-Ck-inversion}
 \Sigma_s(A,B)
 =\sum_{k=s}^r(-1)^{k-s}\binom{k}{s}C_k(A,B),
 \qquad 0\le s\le r.
\end{equation}
\end{proposition}

\noindent The proof is given in
Appendix~\ref{app:proof-margin-inversion}.

Since $\widehat\eta_{n,r,s}=\Sigma_s/N_{n,r,s}$,
Proposition~\ref{prop:margin-inversion} evaluates the direct estimator
\eqref{eq:direct-general} from the $r+1$ quantities $C_0,\ldots,C_r$.
Each $r$-subset contributes to the margins indexed by its $2^r$ subsets, so
all margins can be accumulated in $O(2^r\binom nr)$ additions.  Forming their
cross-products and applying \eqref{eq:Sigma-Ck-inversion} has the same
$O(n^r)$ order for fixed $r$.

For $r=2$, these margin cross-products are related to the terms in
\textup{(\ref{eq:pair-contraction})} by
\[
 C_0(A,B)=\tfrac14A_{\cdot\cdot}B_{\cdot\cdot},\qquad
 C_1(A,B)=\sum_{i=1}^n A_{i\cdot}B_{i\cdot},\qquad
 C_2(A,B)=\tfrac12\sum_{i\ne j}A_{ij}B_{ij}.
\]
Consequently, the right-hand side of \textup{(\ref{eq:pair-contraction})} is
\[
 \frac{1}{n(n-3)}
 \left\{
 2C_2(A,B)-\frac{2}{n-2}C_1(A,B)
 +\frac{4}{(n-1)(n-2)}C_0(A,B)
 \right\},
\]
which recovers the earlier pairwise inner-product formula.


\subsection{Projection representation by higher-order U-centering}
The margin cross-products from the preceding subsection also enter the ANOVA
projector, since
\begin{equation*}
 C_k(A,B)=\inner{A}{D_{k,r}^{\ast}D_{k,r}B}.
\end{equation*}
The result below shows that the appropriate linear combination equals the
cross-product of the projected residual arrays.

Relative to the coordinate basis indexed by $\cI_{n,r}$, invariance under
relabeling implies that the entries of $\cU_{n,r}$ depend only on intersection
size; see \citet{Dunkl1978,Filmus2016}.  For $I,J\in\cI_{n,r}$ with
$|I\cap J|=s$, set
\begin{equation}\label{eq:p-s-def}
 p_{n,r,s}
 =(\cU_{n,r})_{I,J}
 =\sum_{j=0}^s
   \frac{(-1)^{r-j}\binom sj}{\binom{n-r-j+1}{r-j}}.
\end{equation}
The following closed form gives
the coefficient normalization used in Theorem~\ref{thm:general-compression}
to identify the projected inner product in \textup{(\ref{eq:general-compression})}
with the direct overlap estimator in \textup{(\ref{eq:direct-general})}.

\begin{lemma}
\label{lem:p-s}
Let $r\ge1$, $n\ge2r$, and $\kappa=n-2r+1$.  For $0\le s\le r$,
\begin{equation}\label{eq:p-s-closed}
 p_{n,r,s}
 =(-1)^{r-s}
 \frac{\kappa(r-s)!(\kappa+s-1)!}{(\kappa+r)!}
 =(-1)^{r-s}
 \frac{(n-2r+1)(r-s)!(n-2r+s)!}{(n-r+1)!}.
\end{equation}
Moreover,
\begin{equation}\label{eq:group-coeff}
 \frac{N_{n,r,s}p_{n,r,s}}{d_{n,r}}
 =(-1)^{r-s}\binom rs.
\end{equation}
\end{lemma}

The proof is given in Appendix~\ref{app:p-s-proof}.

\begin{samepage}
\begin{theorem}[Projection representation and unbiasedness]
\label{thm:general-compression}
Let $a$ and $b$ be symmetric square-integrable kernels with $r$ arguments,
and let $n\ge2r$.  Form the arrays $A_I=a(X_I)$ and $B_I=b(Y_I)$ on
$\cI_{n,r}$.  Then
\begin{equation}\label{eq:general-compression}
 \boxed{
 T_{n,r}^{\mathrm{dir}}
 =\frac{1}{d_{n,r}}
   \inner{\cU_{n,r}A}{\cU_{n,r}B}
 =\frac{1}{d_{n,r}}\inner{A}{\cU_{n,r}B}
 =\frac{1}{d_{n,r}}\inner{\cU_{n,r}A}{B}.
 }
\end{equation}
In particular,
\begin{equation*}
 \E\left[
 \frac{1}{d_{n,r}}
 \inner{\cU_{n,r}A}{\cU_{n,r}B}
\right]
 =\Theta_r(a,b).
\end{equation*}
\end{theorem}
\end{samepage}


\begin{proof}
Since $\cU_{n,r}$ is self-adjoint and idempotent,
\[
 \inner{\cU_{n,r}A}{\cU_{n,r}B}
 =\inner{A}{\cU_{n,r}B}.
\]
Expanding the latter inner product and grouping the ordered pairs $(I,J)$
according to $s=|I\cap J|$, the definitions of $p_{n,r,s}$,
$\Sigma_s(A,B)$, and $\widehat\eta_{n,r,s}$ give
\begin{align*}
 \frac{1}{d_{n,r}}\inner{A}{\cU_{n,r}B}
 &=\frac{1}{d_{n,r}}\sum_{s=0}^r
   p_{n,r,s}\Sigma_s(A,B)\\
 &=\sum_{s=0}^r
   \frac{N_{n,r,s}p_{n,r,s}}{d_{n,r}}
   \widehat\eta_{n,r,s}\\
 &=\sum_{s=0}^r(-1)^{r-s}\binom rs
   \widehat\eta_{n,r,s}
 =T_{n,r}^{\mathrm{dir}}.
\end{align*}
The final line uses \eqref{eq:group-coeff} and the definition
\eqref{eq:direct-general}.
Unbiasedness follows from Theorem~\ref{thm:overlap-anova}.
\end{proof}

From the raw arrays, accumulating the margins and forming $\cU_{n,r}B$ each
cost $O(2^r\binom nr)$, since every entry touches $2^r$ margins.  The one-sided
identity $\inner{A}{\cU_{n,r}B}/d_{n,r}$ adds only $O(\binom nr)$, giving
$O(2^r\binom nr)=O(n^r)$ arithmetic for fixed $r$.

For the matching kernels in Section~\ref{sec:example}, the sample
calculation uses the observable raw arrays $A_I=m_k(X_I)$ and
$B_I=m_\ell(Y_I)$; the population-centered kernels are needed only to
interpret the target.  When $n\ge8$, a literal implementation of the direct
overlap estimator contains terms involving as many as eight distinct
observations.  Proposition~\ref{prop:margin-inversion} and
Theorem~\ref{thm:general-compression} evaluate exactly the same unbiased
estimator of $H_4$ in $O(n^4)$ arithmetic after the pairwise Gram entries have
been formed.  Although $H_4$ is nonnegative, this unbiased cross-product
estimator can be negative in finite samples.


The projected residual arrays also have zero lower-order margins.  Their
entrywise conditional interpretation is given next.

\begin{corollary}
\label{cor:general-conditional}
Under the assumptions and notation of Theorem~\ref{thm:general-compression},
for a fixed $I\in\cI_{n,r}$,
\begin{equation}\label{eq:general-conditional}
 \E\left[(\cU_{n,r}A)_I\given X_I\right]
 =\lambda_{n,r}a_r(X_I),
 \qquad
 \lambda_{n,r}
 =\frac{d_{n,r}}{\binom nr}
 =\frac{n-2r+1}{n-r+1}.
\end{equation}
\end{corollary}

\noindent The proof is given in
Appendix~\ref{app:proof-general-conditional}.


\subsection{\texorpdfstring{{The explicit estimator for three-argument kernels}}{The explicit estimator for three-argument kernels}}

We now specialize the general estimator to $r=3$.  Assume $n\ge6$.  For
symmetric three-argument kernels $a$ and $b$, define the arrays
$A_{ijk}=a(X_i,X_j,X_k)$ and $B_{ijk}=b(Y_i,Y_j,Y_k)$ on distinct triples.
The target is
\[
 \Theta_3(a,b)=\E\{a_3(X_1,X_2,X_3)b_3(Y_1,Y_2,Y_3)\}.
\]
The direct estimator is
\begin{equation*}
 \widehat\Theta_{n,3}
 =\widehat\eta_{n,3,3}
 -3\widehat\eta_{n,3,2}
 +3\widehat\eta_{n,3,1}
 -\widehat\eta_{n,3,0},
\end{equation*}
which involves up to six distinct observations.  By
Theorem~\ref{thm:general-compression},
\begin{equation*}
 \widehat\Theta_{n,3}
 =\frac{1}{\binom n3-\binom n2}
   \sum_{i<j<k}
   (\cU_{n,3}A)_{ijk}(\cU_{n,3}B)_{ijk},
\end{equation*}
where each centered array is given by \eqref{eq:r3-formula}.  After the
arrays for the three-argument kernels have been formed, the exact centered
calculation uses $O(n^3)$ arithmetic operations instead of the $O(n^6)$
products in a literal enumeration over pairs of triples.


\subsection{Application: residual representation for the variance of the
highest-order Hoeffding component}
\label{subsec:variance-components}

Assume $n\ge 2r$.  Let $Z_1,\ldots,Z_n$ be i.i.d.\ copies of a random variable
$Z$ taking values in $\mathcal Z$.  For $I=\{i_1<\cdots<i_r\}$, write
$Z_I=(Z_{i_1},\ldots,Z_{i_r})$.  Let $h:\mathcal Z^r\to\R$ be a symmetric
square-integrable kernel, and define
\begin{equation*}
 U_n(h)=\binom nr^{-1}\sum_{I\in\cI_{n,r}}h(Z_I).
\end{equation*}
Write $h_1,\ldots,h_r$ for its Hoeffding components and
\begin{equation*}
 \sigma_j^2=\E\{h_j(Z_1,\ldots,Z_j)^2\}.
\end{equation*}
Orthogonality of the Hoeffding decomposition gives the exact finite-sample
variance formula
\begin{equation*}
 \Var\{U_n(h)\}
 =\sum_{j=1}^r\frac{\binom rj^2}{\binom nj}\sigma_j^2.
\end{equation*}
This variance decomposition is classical; see
\citet{Hoeffding1948,Serfling1980}.  Here we focus only on $\sigma_r^2$, the
variance of the highest-order Hoeffding component.

\begin{samepage}
Let $H\in\cV_{n,r}$ be the kernel array with $H_I=h(Z_I)$.  Taking $A=B=H$
in Theorem~\ref{thm:general-compression} identifies the corresponding
overlap-based estimator with the normalized squared residual norm
\begin{equation*}
 \widehat\sigma_{r,n}^2
 :=T_{n,r}^{\mathrm{dir}}
 =\frac{1}{d_{n,r}}\norm{\cU_{n,r}H}^2.
\end{equation*}
Its expectation is $\Theta_r(h,h)=\sigma_r^2$, and it is nonnegative.
\end{samepage}

For an order-three U-statistic,
\begin{equation*}
 \Var\{U_n(h)\}
 =\frac{9}{n}\sigma_1^2
  +\frac{9}{\binom n2}\sigma_2^2
  +\frac{1}{\binom n3}\sigma_3^2,
\end{equation*}
and
\begin{equation*}
 \widehat\sigma_{3,n}^2
 =\frac{1}{\binom n3-\binom n2}
   \sum_{i<j<k}(\cU_{n,3}H)_{ijk}^2.
\end{equation*}
If $h_1=h_2=0$---equivalently, if $h-\E h$ is canonical of order three---then
$\Var\{U_n(h)\}=\sigma_3^2/\binom n3$, so
$\widehat\sigma_{3,n}^2/\binom n3$ is a nonnegative unbiased variance
estimator.



\section{Discussion}
\label{sec:discussion}
\looseness=-1
U-centering is least-squares residualization on sample subsets, and the
normalized residual cross-product equals the direct overlap estimator of the
$r$th Hoeffding pairing.  This quantity can be evaluated in $O(n^r)$
operations for fixed $r$.  When both arrays are formed from the same kernel
and sample, the projection identity identifies the overlap-based estimator of
the variance of the highest-order Hoeffding component with a nonnegative
squared residual norm.  Thus the familiar pairwise formula is the
$r=2$ case of a broader bilinear construction; general higher-order
U-statistics outside this form require different computational methods.

At the pairwise level, U-centering turns two familiar test
statistics into regression quantities.  The U-centered dependence
$t$-statistic is the slope $t$-statistic after additive endpoint effects have
been removed.  In the two-sample problem, $\mathcal E_{n,m}/2$ is the slope on
the between-group pair indicator.  The common-endpoint and fully interacted
regressions estimate the same slope, but the fully interacted model adds
$N-2=n+m-2$ endpoint-by-pair-type effects.  These extra effects allow an
observation's endpoint effect to differ between within- and between-group
pairs.  This regression representation explains the pooled denominator in
\citet{ChakrabortyZhang2021} and the alternative restricted denominator.  The
same regression interpretation extends to arrays indexed by
higher-order subsets after all lower-order subset effects have been removed;
see Corollary~\ref{cor:higher-order-subset-regression}.

\appendix

\section{The higher-order ANOVA projection}
\label{app:higher-projection-proof}
\begin{proof}
\smallskip
\noindent\emph{Self-adjointness and zero margins.}
Each operator $D_{j,r}^\ast D_{j,r}$ is self-adjoint, so
$\cU_{n,r}$ is self-adjoint.  We now compute its $(r-1)$-margins.

Fix $S\in\binom{[n]}{r-1}$.  For $0\le j\le r$, write
\[
 m_j(J)=(D_{j,r}A)(J),
 \qquad
 T_j(S)=\sum_{\substack{J\subseteq S\\|J|=j}}m_j(J),
\]
with $T_{-1}(S)=0$.  The contribution of
$D_{j,r}^\ast D_{j,r}A$ to the $S$-margin is
\begin{equation}\label{eq:Lj-def}
 L_j(S)
 :=\sum_{i\notin S}
   \sum_{\substack{J\subseteq S\cup\{i\}\\|J|=j}}m_j(J).
\end{equation}
Split the inner $j$-sets according to whether they contain the new label $i$.
Those not containing $i$ are the $j$-subsets of $S$; each is repeated for all
$n-r+1$ choices of $i\notin S$.  Their total contribution is
$(n-r+1)T_j(S)$.

For $j=0$, there is no second part.  For $j\ge1$, a $j$-set containing
$i$ has the form $K\cup\{i\}$ with $K\subseteq S$ and $|K|=j-1$.
For a fixed such $K$, we evaluate the sum over labels outside $S$.  Since
$K\subseteq S$, the complement of $K$ has the disjoint decomposition
\[
 [n]\setminus K
 =\bigl([n]\setminus S\bigr)\,\dot\cup\,\bigl(S\setminus K\bigr).
\]
Consequently,
\begingroup
\interdisplaylinepenalty=10000
\begin{align}
 \sum_{i\notin S}m_j(K\cup\{i\})
 &=\sum_{i\notin K}m_j(K\cup\{i\})
   -\sum_{u\in S\setminus K}m_j(K\cup\{u\}) \notag\\
 &=\sum_{i\notin K}
   \sum_{\substack{I\in\cI_{n,r}\\I\supseteq K\cup\{i\}}}A_I
   -\sum_{u\in S\setminus K}m_j(K\cup\{u\}) \notag\\
 &=\sum_{\substack{I\in\cI_{n,r}\\I\supseteq K}}
   \sum_{i\in I\setminus K}A_I
   -\sum_{u\in S\setminus K}m_j(K\cup\{u\}) \notag\\
 &=(r-j+1)m_{j-1}(K)
   -\sum_{u\in S\setminus K}m_j(K\cup\{u\}).
 \label{eq:fixed-K-count}
\end{align}
\endgroup
The first equality separates the labels outside $K$ into those outside $S$
and those in $S\setminus K$.  The second equality expands the definition of
$m_j$, and the third reverses the order of summation.  For every fixed
$r$-set $I\supseteq K$, the inner sum contains one copy of $A_I$ for each
$i\in I\setminus K$.  Since $|K|=j-1$,
\[
 |I\setminus K|=r-(j-1)=r-j+1,
\]
so each $A_I$ is counted exactly $r-j+1$ times.  This gives
$(r-j+1)m_{j-1}(K)$.  Finally, subtracting the terms indexed by
$u\in S\setminus K$ removes precisely those labels that are outside $K$ but
still inside $S$, leaving only the labels $i\notin S$ required in
\eqref{eq:Lj-def}.

Next, sum \eqref{eq:fixed-K-count} over all $(j-1)$-subsets $K$ of $S$.
The first term becomes $(r-j+1)T_{j-1}(S)$.  In the final double sum, each
$j$-subset $J\subseteq S$ appears once for each choice of the removed element
$u\in J$, hence exactly $j$ times.  Combining the two parts yields
  \begin{equation*}
 L_j(S)
 =(n-r+1-j)T_j(S)+(r-j+1)T_{j-1}(S).
\end{equation*}

Because
\[
 D_{r-1,r}(\cU_{n,r}A)(S)=\sum_{j=0}^r c_{n,r,j}L_j(S),
\]
the coefficient of $T_j(S)$, for $0\le j\le r-1$, is
\[
 (n-r+1-j)c_{n,r,j}+(r-j)c_{n,r,j+1}.
\]
Setting these coefficients to zero yields
\[
 c_{n,r,j+1}
 =-\frac{n-r-j+1}{r-j}c_{n,r,j},
 \qquad 0\le j<r.
\]
This is the recurrence stated before Definition~\ref{def:higher-u}; the
coefficients used there satisfy it with $c_{n,r,r}=1$.  Since $T_r(S)=0$
for $|S|=r-1$, all terms cancel.  Therefore
$D_{r-1,r}(\cU_{n,r}A)=0$, proving that the range of $\cU_{n,r}$ is
contained in $\cH_{n,r}$.

\medskip
\noindent\emph{Identity on the residual space.}
Conversely, suppose $A\in\cH_{n,r}$.  For a $j$-set $J$ with $j<r-1$,
sum the zero $(r-1)$-margins over all $(r-1)$-sets $K$ containing $J$:
\[
 0=\sum_{\substack{K\supseteq J\\|K|=r-1}}
   (D_{r-1,r}A)(K)
  =(r-j)(D_{j,r}A)(J).
\]
Each $r$-set containing $J$ is counted $r-j$ times, so every lower margin
$D_{j,r}A$ vanishes.  In \eqref{eq:higher-u-operator}, only the $j=r$ term
remains, and that term is $A$.  Hence $\cU_{n,r}A=A$ on $\cH_{n,r}$.

\medskip
\noindent\emph{Projection and kernel.}
A self-adjoint map whose range lies in $\cH_{n,r}$ and which is the identity on
$\cH_{n,r}$ is the orthogonal projection onto that space.  Its kernel is
therefore $\cH_{n,r}^{\perp}=\cL_{n,r}$.

\medskip
\noindent\emph{Dimension.}
The inclusion matrix $D_{r-1,r}$ has full row rank
$\binom n{r-1}$ when $n\ge2r-1$; this is the classical rank theorem of
\citet{Gottlieb1966}.  Therefore
\[
 \dim\cH_{n,r}
 =\dim\ker(D_{r-1,r})
 =\binom nr-\binom n{r-1},
\]
which is \eqref{eq:general-df}.
\end{proof}


\section{\texorpdfstring{{RKHS identities for four-argument matching kernels}}{RKHS identities for four-argument matching kernels}}
\label{app:rkhs}
\noindent\emph{Centered feature maps.}
Let $k$ and $\ell$ be the bounded measurable positive-definite kernels from Section~\ref{sec:example}, with separable RKHSs $\mathcal F$ and
$\mathcal G$.  Their canonical feature maps are
\[
 \phi(x)=k(x,\cdot)\in\mathcal F,
 \qquad
 \psi(y)=\ell(y,\cdot)\in\mathcal G.
\]
The reproducing property gives
\[
 k(x,x')=\inner{\phi(x)}{\phi(x')}_{\mathcal F},
 \qquad
 \ell(y,y')=\inner{\psi(y)}{\psi(y')}_{\mathcal G}.
\]
Boundedness ensures that the mean embeddings
$\mu_X=\E\phi(X)$ and $\mu_Y=\E\psi(Y)$ exist.  With
$\phi_c=\phi-\mu_X$ and $\psi_c=\psi-\mu_Y$, the population-centered kernels
satisfy
\[
 k_c(x,x')=\inner{\phi_c(x)}{\phi_c(x')}_{\mathcal F},
 \qquad
 \ell_c(y,y')=\inner{\psi_c(y)}{\psi_c(y')}_{\mathcal G}.
\]

\medskip
\noindent\emph{Fourth Hoeffding projections.}
In each matching term, the two pair kernels involve disjoint arguments, so
the full four-coordinate Hoeffding projection replaces each raw pair kernel
by its second Hoeffding component.  Let $\mathbb E_i$ denote integration in
the $i$th argument.  For example,
\[
 \left\{\prod_{i=1}^4(I-\mathbb E_i)\right\}
       \{k_{12}k_{34}\}
 =\{(I-\mathbb E_1)(I-\mathbb E_2)k_{12}\}
  \{(I-\mathbb E_3)(I-\mathbb E_4)k_{34}\}
 =k_{c,12}k_{c,34}.
\]
Integrating out any one argument annihilates the unique centered factor
containing that argument, which also verifies canonicality.  Relabeling the
arguments and averaging over the three matchings gives the two displayed
fourth-component formulas in Section~\ref{sec:example}.

\medskip
\noindent\emph{Cross-covariance operator.}
Recall from Section~\ref{sec:example} that the cross-covariance operator is
\[
 C_{XY}=\E\{\phi_c(X)\otimes\psi_c(Y)\}.
\]
It is well defined in the Hilbert--Schmidt class: boundedness makes the centered
features uniformly bounded, while separability supplies the required
measurability.  Consequently, $C_{XY}^{\ast}C_{XY}$ is positive and
trace-class.  The operator also satisfies, for $f\in\mathcal F$ and
$g\in\mathcal G$,
\[
 \inner{f}{C_{XY}g}_{\mathcal F}
 =\E\!\left[
   \inner{f}{\phi_c(X)}_{\mathcal F}
   \inner{\psi_c(Y)}{g}_{\mathcal G}
 \right]
 =\Cov\{f(X),g(Y)\}.
\]

\medskip
\noindent\emph{Traces and spectral identities.}
For a positive trace-class operator $T$ on $\mathcal G$, its trace is
\[
 \tr(T)=\sum_q\inner{f_q}{Tf_q}_{\mathcal G},
\]
where the value is independent of the orthonormal basis $(f_q)$.
The definitions in Section~\ref{sec:example} give
\[
 \norm{C_{XY}}_{\mathrm{HS}}^2
 =\tr(C_{XY}^{\ast}C_{XY})
 =\sum_j s_j^2,
 \qquad
 \tr\{(C_{XY}^{\ast}C_{XY})^2\}=\sum_j s_j^4.
\]
Finally, for independent copies $(X_1,Y_1)$ and $(X_2,Y_2)$ of $(X,Y)$,
\[
 \norm{C_{XY}}_{\mathrm{HS}}^2
 =\E\{k_c(X_1,X_2)\ell_c(Y_1,Y_2)\}.
\]
The two coordinates within each pair may remain dependent.  These identities
justify \eqref{eq:H2-spectrum}.

\medskip
\noindent\emph{Matching expansion for \eqref{eq:H4-spectrum}.}
For the representative four-cycle term in Section~\ref{sec:example}, choose
orthonormal bases $(e_p)$ and $(f_q)$ of $\mathcal F$ and $\mathcal G$, and
set
\[
 c_{pq}
 =\E\!\left[
   \inner{\phi_c(X)}{e_p}
 \inner{\psi_c(Y)}{f_q}
 \right].
\]
Put
\[
 \xi_{ip}=\inner{\phi_c(X_i)}{e_p},
 \qquad
 \zeta_{iq}=\inner{\psi_c(Y_i)}{f_q}.
\]
Parseval's identity gives
\[
 k_{c,ij}=\sum_p\xi_{ip}\xi_{jp},
 \qquad
 \ell_{c,ij}=\sum_q\zeta_{iq}\zeta_{jq}.
\]
For finite truncations of the basis expansions, the calculation below is
elementary.  Each truncated kernel expansion is bounded by the corresponding
product of centered feature norms, and these norms are uniformly bounded
because the kernels are bounded.  Dominated convergence therefore permits
passage to the infinite sums.  Using independence across the four
paired-observation labels, we obtain
\begin{align*}
 &\E\{k_{c,12}k_{c,34}\ell_{c,13}\ell_{c,24}\}\\
 &\quad=
 \sum_{p,p',q,q'}
 \E(\xi_{1p}\zeta_{1q})
 \E(\xi_{2p}\zeta_{2q'})
 \E(\xi_{3p'}\zeta_{3q})
 \E(\xi_{4p'}\zeta_{4q'})\\
 &\quad=
 \sum_{p,p',q,q'}
 c_{pq}c_{pq'}c_{p'q}c_{p'q'}.
\end{align*}
Here dependence between $X_i$ and $Y_i$ within a pair is retained.  Since
$c_{pq}=\inner{C_{XY}f_q}{e_p}$, let
\[
 T:=C_{XY}^{\ast}C_{XY}:\mathcal G\to\mathcal G.
\]
This is a positive self-adjoint operator.  Its coefficients relative to the
orthonormal basis $(f_q)$ are
\[
\begin{aligned}
 t_{qq'}
 &:=\inner{Tf_{q'}}{f_q}_{\mathcal G}
 =\inner{C_{XY}f_{q'}}{C_{XY}f_q}_{\mathcal F}\\
 &=\sum_p c_{pq'}c_{pq}.
\end{aligned}
\]
Therefore
\[
 \sum_{p,p',q,q'}c_{pq}c_{pq'}c_{p'q}c_{p'q'}
 =\sum_{q,q'}t_{qq'}t_{q'q}
 =\tr(T^2)
 =\sum_{j\ge1}s_j^4.
\]
The other five four-cycle terms agree by relabeling, establishing the
identity used in \eqref{eq:H4-spectrum}.


\section{Direct proofs for the pairwise case}
\label{app:pairwise-proofs}

\subsection{Normal equations, range, and null space
  (Corollary~\ref{cor:pair-basic})}
\label{app:proof-pair-basic}
\begin{proof}
Edge vectorization is an isometry from
$\cV_{n,2}$, equipped with $\inner{\cdot}{\cdot}_E$, to
$\R^{\binom{n}{2}}$.  For
$C\in\cV_{n,2}$, let $c=\operatorname{vec}_E(C)$.  Then
\[
 (Z^\mathsf Tc)_i=\sum_{j\ne i}C_{ij}.
\]
Thus $\ker(Z^\mathsf T)$ corresponds exactly to $\cH_{n,2}$.  Likewise,
$\operatorname{col}(Z)$ corresponds exactly to $\cL_{n,2}$ because
$(Zu)_{\{i,j\}}=u_i+u_j$.  Since $Z^\mathsf TZ$ is nonsingular, the matrix
$M_{n,2}$ in \eqref{eq:pair-projection} is the orthogonal projector onto
$\ker(Z^\mathsf T)$ and has kernel $\operatorname{col}(Z)$.  Translating back
to arrays gives
\[
 \cV_{n,2}=\cL_{n,2}\oplus\cH_{n,2},\qquad
 \range(\cU_{n,2})=\cH_{n,2},\qquad
 \ker(\cU_{n,2})=\cL_{n,2}.
\]
An orthogonal projector is self-adjoint and idempotent.  Every element of
$\cH_{n,2}$ has zero row sums, and symmetry gives zero column sums.  Finally,
the orthogonal decomposition of $A$ into its $\cL_{n,2}$ and $\cH_{n,2}$
components proves the stated existence and uniqueness characterization.
\end{proof}

\subsection{Pairwise overlap identity
  (Proposition~\ref{prop:pair-overlap})}
\label{app:proof-pair-overlap}
\begin{proof}
Expand both raw kernels using \eqref{eq:pair-hoeffding}.  In any product term,
if a sample label appears in a nonconstant Hoeffding component on only one
side, integrate over that label while holding the others fixed; canonicality
makes the expectation zero.  Consequently, with no shared label only the
overall means remain.  With one shared label, the pairing of the first
Hoeffding components also remains, and with two shared labels, the pairing of
the second components remains as well.  These three cases give the stated
identities and their weighted difference \eqref{eq:theta2-overlap}.
\end{proof}

\subsection{Conditional interaction interpretation
  (Proposition~\ref{prop:pair-conditional})}
\label{app:proof-pair-conditional}
\begin{proof}
Let $m_a(x)=\E\{a(x,X)\}$.  Conditional on $X_i,X_j$,
\[
 \E(A_{i\cdot}\mid X_i,X_j)
 =a(X_i,X_j)+(n-2)m_a(X_i),
\]
and similarly for $A_{j\cdot}$.  The ordered grand sum has conditional mean
\[
 2a(X_i,X_j)+2(n-2)\{m_a(X_i)+m_a(X_j)\}
 +(n-2)(n-3)a_0.
\]
Substitution into \eqref{eq:pairwise-u} gives
\eqref{eq:pair-conditional}.
\end{proof}


\section{\texorpdfstring{{Auxiliary proofs for the general $r$-subset theory}}{Auxiliary proofs for the general r-subset theory}}
\label{app:general-proofs}

\subsection{Overlap sums from margin cross-products
  (Proposition~\ref{prop:margin-inversion})}
\label{app:proof-margin-inversion}
\begin{proof}
Expanding the two margins and reversing the order of summation gives
\begin{align*}
 C_k(A,B)
 &=\sum_{K\in\binom{[n]}k}
   \sum_{I\supseteq K}\sum_{J\supseteq K}A_I B_J\\
 &=\sum_{I,J\in\cI_{n,r}}
   \#\{K\in\tbinom{[n]}k:K\subseteq I\cap J\}\,A_I B_J\\
 &=\sum_{I,J\in\cI_{n,r}}\binom{|I\cap J|}{k}A_I B_J,
\end{align*}
which is \eqref{eq:Ck-Sigma}.  Equation
\eqref{eq:Sigma-Ck-inversion} follows from the binomial-inversion identity
\[
 \sum_{k=s}^t(-1)^{k-s}\binom{k}{s}\binom{t}{k}
 =\mathbf 1\{t=s\}.
\]
\end{proof}

\subsection{Closed form for the projection entries
  (Lemma~\ref{lem:p-s})}
\label{app:p-s-proof}
\begin{proof}
To prove \eqref{eq:p-s-closed}, recall that $\kappa=n-2r+1$ and write
$q=s-j$ in \eqref{eq:p-s-def} to obtain
\[
 p_{n,r,s}
 =(-1)^{r-s}\sum_{q=0}^s
 \frac{(-1)^q\binom{s}{q}}
      {\binom{\kappa+r-s+q}{r-s+q}}.
\]
For every integer $m\ge0$,
\[
\frac{1}{\binom{\kappa+m}{m}}
 =\kappa\int_0^1t^m(1-t)^{\kappa-1}\,dt,
\]
\begin{samepage}
so
\begin{align*}
 p_{n,r,s}
 &=(-1)^{r-s}\kappa\int_0^1
   t^{r-s}(1-t)^{\kappa-1}
   \sum_{q=0}^s(-1)^q\binom sqt^q\,dt\\
 &=(-1)^{r-s}\kappa\int_0^1
   t^{r-s}(1-t)^{\kappa+s-1}\,dt,
\end{align*}
\end{samepage}
The beta-integral identity gives
\[
 \kappa\int_0^1 t^{r-s}(1-t)^{\kappa+s-1}\,dt
 =\kappa B(r-s+1,\kappa+s)
 =\frac{\kappa(r-s)!(\kappa+s-1)!}{(\kappa+r)!},
\]
which proves \eqref{eq:p-s-closed}.  For the normalization identity
\eqref{eq:group-coeff}, note that
\[
 d_{n,r}=\binom nr\frac{\kappa}{n-r+1}.
\]
Substituting \eqref{eq:N-s} and \eqref{eq:p-s-closed}, and using
\[
 \binom{n-r}{r-s}(r-s)!
 =\frac{(n-r)!}{(n-2r+s)!},
\]
gives \eqref{eq:group-coeff} after cancellation.
\end{proof}

\subsection{\texorpdfstring{{Conditional recovery of the $r$th Hoeffding component
  (Corollary~\ref{cor:general-conditional})}}{Conditional recovery of the rth Hoeffding component}}
\label{app:proof-general-conditional}
\begin{proof}
Decompose the random array using \eqref{eq:general-hoeffding-decomp}.  For each
$j<r$, the $j$th-order array contribution
\[
 I\longmapsto
 \sum_{\substack{S\subseteq I\\|S|=j}}a_j(X_S)
\]
belongs to $\cL_{n,r}$ and is therefore annihilated by $\cU_{n,r}$.  Write
$A_I^{(r)}=a_r(X_I)$.  Then
\[
 \cU_{n,r}A=\cU_{n,r}A^{(r)}.
\]
Conditional on $X_I$, the canonicality of $a_r$ implies
$\E\{a_r(X_J)\mid X_I\}=0$ whenever $J\ne I$.  Therefore only the diagonal
projection entry remains.  By permutation symmetry, all
$\binom nr$ diagonal entries of the matrix representing $\cU_{n,r}$ are
equal.  Since $\cU_{n,r}$ is an orthogonal projection onto the
$d_{n,r}$-dimensional space $\cH_{n,r}$ by
Theorem~\ref{thm:higher-projection}, the sum of its diagonal entries is
$d_{n,r}$; hence each diagonal entry is
$d_{n,r}/\binom nr$.  This proves \eqref{eq:general-conditional}.
\end{proof}


\section{Algebraic proofs for the two-sample edge regressions}
\label{app:two-sample-nesting}

We retain the notation of Section~\ref{sec:edge-regression}.
The dependence regression follows from
Proposition~\ref{prop:general-edge-regression}, with $W=Z$, together with
Theorem~\ref{thm:pair-anova}.  We therefore focus on
Propositions~\ref{prop:energy-group-slope} and~\ref{prop:energy-anova-t}.
We first project the between-group pair indicator in the common-endpoint
model.  We then analyze the fully interacted model block by block, recover
its pooled studentization, and compare the two fitted spaces.

\noindent\emph{Projection of the between-group indicator and the common
slope.}
Recall that $M_{N,2}$ projects onto $\ker(Z_N^\mathsf T)$, the space
orthogonal to all additive endpoint effects.  To prove
\eqref{eq:ucentered-group-indicator}, we show that
$\ell_E/a_{n,m}^2$ belongs to this residual space and that its difference
from $c_E$ belongs to $\operatorname{col}(Z_N)$.

The $i$th entry of $Z_N^\mathsf T\ell_E$ is the sum of the entries of
$\ell_E$ over the edges incident to label $i$.  For an $X$-sample label and
a $Y$-sample label, respectively, these sums are
\[
 (n-1)\left\{-\frac1{n(n-1)}\right\}
 +m\left(\frac1{nm}\right)=0,
 \qquad
 (m-1)\left\{-\frac1{m(m-1)}\right\}
 +n\left(\frac1{nm}\right)=0.
\]
Hence $Z_N^\mathsf T\ell_E=0$, and therefore
$\ell_E\in\ker(Z_N^\mathsf T)=\range(M_{N,2})$.  Moreover, only the $nm$
between-group edges contribute to $c_E^\mathsf T\ell_E$, so
\[
 c_E^\mathsf T\ell_E=nm\left(\frac1{nm}\right)=1,
 \qquad
 \norm{\ell_E}^2=a_{n,m}^2.
\]

It remains to show that the difference between $c_E$ and
$\ell_E/a_{n,m}^2$ is an additive endpoint effect.  Set
\[
 u_X=\frac{1}{2n(n-1)a_{n,m}^2},
 \qquad
 u_Y=\frac{1}{2m(m-1)a_{n,m}^2}.
\]
On an $XX$ edge, the entry of $c_E-\ell_E/a_{n,m}^2$ is $2u_X$, and on a
$YY$ edge it is $2u_Y$.  On an $XY$ edge, the definition of $a_{n,m}^2$
gives
\[
 1-\frac{1}{nm a_{n,m}^2}=u_X+u_Y.
\]
Let $u\in\R^N$ equal $u_X$ on the $X$-sample labels and $u_Y$ on the
$Y$-sample labels.  For an edge $\{i,j\}$, the corresponding entry of
$Z_Nu$ is $u_i+u_j$.  The three edge-type calculations above therefore give
\[
 c_E-\frac{\ell_E}{a_{n,m}^2}=Z_Nu.
\]
This is the orthogonal decomposition of $c_E$ into a vector in
$\range(M_{N,2})$ and a vector in
$\ker(M_{N,2})=\operatorname{col}(Z_N)$.  Applying $M_{N,2}$ proves
\eqref{eq:ucentered-group-indicator}.

Using $M_{N,2}c_E=\ell_E/a_{n,m}^2$,
$c_E^\mathsf T\ell_E=1$, and the symmetry of $M_{N,2}$, the
Frisch--Waugh--Lovell formula gives
\[
 \widehat\theta
 =\frac{c_E^\mathsf TM_{N,2}a_E}
        {c_E^\mathsf TM_{N,2}c_E}
 =\frac{\ell_E^\mathsf Ta_E/a_{n,m}^2}
        {1/a_{n,m}^2}
 =\ell_E^\mathsf Ta_E
 =\overline A_{XY}-\frac{\overline A_{XX}+\overline A_{YY}}2
 =\frac{\mathcal E_{n,m}}2.
\]
The penultimate equality follows by grouping the entries of
$\ell_E^\mathsf Ta_E$ over the $XX$, $XY$, and $YY$ blocks.

By Theorem~\ref{thm:pair-anova}, the edge vectors of $\widetilde A$ and
$\widetilde C$ are $M_{N,2}a_E$ and $M_{N,2}c_E$, respectively.  Because
$M_{N,2}$ is an orthogonal projector,
\[
 \inner{\widetilde A}{\widetilde C}_E
 =(M_{N,2}a_E)^\mathsf T(M_{N,2}c_E)
 =\frac{\ell_E^\mathsf Ta_E}{a_{n,m}^2}
 =\frac{\mathcal E_{n,m}}{2a_{n,m}^2},
\]
which proves \eqref{eq:energy-as-ucentered-label-covariance}.

The same calculation also gives the restricted-model residual sum of
squares.  Adding the residualized predictor $M_{N,2}c_E$ reduces the
endpoint-only residual sum of squares by
\[
 \frac{(c_E^\mathsf TM_{N,2}a_E)^2}
      {c_E^\mathsf TM_{N,2}c_E}
 =\frac{\widehat\theta^2}{a_{n,m}^2}.
\]
This proves the formula for $\mathrm{RSS}_R$ in
\eqref{eq:restricted-two-sample-rss}.  The matrix $Z_N$ has rank $N$, and
$M_{N,2}c_E=\ell_E/a_{n,m}^2\ne0$, so adding $c_E$ raises the fitted rank to
$N+1$.  The restricted residual degrees of freedom are therefore
$\binom N2-N-1$, as stated in \eqref{eq:restricted-two-sample-rss}.  Moreover,
$(c_E^\mathsf TM_{N,2}c_E)^{-1}=a_{n,m}^2$ is the least-squares variance
factor for $\widehat\theta$.  Provided $\mathrm{RSS}_R>0$, its usual
standard error is therefore
$a_{n,m}\sqrt{\mathrm{RSS}_R/\nu_R}$, which gives
\eqref{eq:restricted-energy-t}.

\smallskip
\noindent\emph{Blockwise solution of the fully interacted model.}
Write the three block means as
\[
 \mu_{XX}=\eta+\delta,
 \qquad
 \mu_{XY}=\eta+\theta,
 \qquad
 \mu_{YY}=\eta-\delta.
\]
This is a one-to-one reparametrization because
\[
 \eta=\frac{\mu_{XX}+\mu_{YY}}2,
 \qquad
 \delta=\frac{\mu_{XX}-\mu_{YY}}2,
 \qquad
 \theta=\mu_{XY}-\frac{\mu_{XX}+\mu_{YY}}2.
\]
The $X$- and $Y$-group $\alpha$ effects occur only in the $XX$ and $YY$
blocks, respectively, while the $\beta$ effects occur only as row and column
effects in the $XY$ block.  Because the three blocks contain disjoint pair
entries, the full least-squares criterion is a sum of three criteria that can
be minimized separately.

Each centered endpoint-effect term has average zero in its block.  Indeed,
every $\alpha_i$ appears in $n-1$ edges of the $XX$ block, so the sum of the
endpoint terms over that block is
$(n-1)\sum_{i\le n}\alpha_i=0$.  The same argument applies to the $YY$ block.
In the $XY$ block, the corresponding sum is
$m\sum_{i\le n}\beta_i+n\sum_{j>n}\beta_j=0$.
The fitted block means are therefore
$\overline A_{XX}$, $\overline A_{XY}$, and $\overline A_{YY}$.  Substituting
them into the expression for $\theta$ gives
\[
 \widehat\theta
 =\overline A_{XY}
  -\frac{\overline A_{XX}+\overline A_{YY}}2
 =\ell_E^\mathsf Ta_E
 =\frac{\mathcal E_{n,m}}2.
\]

For the rectangular $XY$ block, fitting a mean together with centered row
and column effects is ordinary two-way additive ANOVA.  Double-centering
subtracts the fitted row and column effects, with the grand mean restored.
For the symmetric $XX$ and $YY$ blocks, fitting a mean and additive endpoint
effects is the pairwise ANOVA problem in Theorem~\ref{thm:pair-anova}.  The
endpoint incidence space already contains the constant vector, so
U-centering removes both the block mean and the endpoint effects.  The three
residuals are thus
\[
 \Pi_nA^{XY}\Pi_m,
 \qquad
 M_{n,2}\operatorname{vec}_E(A^{XX}),
 \qquad
 M_{m,2}\operatorname{vec}_E(A^{YY}).
\]
The first residual is a rectangular matrix whose squared norm is the
Frobenius norm; the other two are upper-triangular edge vectors.  Their
squared norms therefore add to \eqref{eq:stacked-full-rss} without any
additional factor of two.

The fitted dimensions of the $XX$, $XY$, and $YY$ blocks are, respectively,
$n$, $N-1$, and $m$.  Hence the corresponding residual spaces have
dimensions
\[
 d_{n,2}=\binom n2-n,
 \qquad
 (n-1)(m-1),
 \qquad
 d_{m,2}=\binom m2-m.
\]
Their sum is
\[
 \nu_F=(n-1)(m-1)+d_{n,2}+d_{m,2}
 =\binom N2-2N+1,
\]
which proves \eqref{eq:stacked-two-sample-df}.  Equivalently, the full fitted
space has rank $n+(N-1)+m=2N-1$.

\smallskip
\noindent\emph{Recovery of the pooled studentization.}
For the same dissimilarity convention as in
\citet{ChakrabortyZhang2021}, let $v_s=s(s-3)/2=d_{s,2}$, and let
$\widetilde A^{XX}$ and $\widetilde A^{YY}$ denote the U-centered
within-group arrays.  In their notation, the cross-block quantity and the two
within-sample squared distance variances are
\begin{align*}
 \operatorname{cdCov}_{n,m}^2(X,Y)
 &=\frac{\norm{\Pi_nA^{XY}\Pi_m}_F^2}{(n-1)(m-1)}
   =\frac{\mathrm{RSS}_{XY}}{(n-1)(m-1)},\\
 \widetilde{\mathcal D}_n^2(X,X)
 &=\frac{1}{n(n-3)}\sum_{i\ne j}(\widetilde A^{XX}_{ij})^2
   =\frac{\mathrm{RSS}_{XX}}{v_n},\\
 \widetilde{\mathcal D}_m^2(Y,Y)
 &=\frac{1}{m(m-3)}\sum_{i\ne j}(\widetilde A^{YY}_{ij})^2
   =\frac{\mathrm{RSS}_{YY}}{v_m}.
\end{align*}
The rectangular $XY$ array lists each cross edge once.  By contrast, the two
within-sample definitions sum over ordered pairs, so symmetry counts every
within-group edge twice.  Substituting the three identities above into the
pooled variance estimator of \citet{ChakrabortyZhang2021} gives
\begin{equation}
 S_{n,m}
 =\frac{4(n-1)(m-1)\operatorname{cdCov}_{n,m}^2(X,Y)
          +4v_n\widetilde{\mathcal D}_n^2(X,X)
          +4v_m\widetilde{\mathcal D}_m^2(Y,Y)}
         {(n-1)(m-1)+v_n+v_m}.
 \label{eq:cz-pooled-variance}
\end{equation}
The three identities show that its numerator is
$4(\mathrm{RSS}_{XY}+\mathrm{RSS}_{XX}+\mathrm{RSS}_{YY})$ and its
denominator is $\nu_F$.  Hence $S_{n,m}=4\mathrm{MS}_F$.
The block-mean calculation above writes
$\widehat\theta=\ell_E^\mathsf Ta_E$, so $\ell_E$ is its least-squares
coefficient-weight vector.  Since $\norm{\ell_E}^2=a_{n,m}^2$, the usual
homoskedastic least-squares formula gives, provided $\mathrm{MS}_F>0$,
\[
 \widehat{\operatorname{se}}(\widehat\theta)
 =a_{n,m}\sqrt{\mathrm{MS}_F}.
\]
Consequently,
\[
 \widehat{\operatorname{se}}(\mathcal E_{n,m})
 =2a_{n,m}\sqrt{\mathrm{MS}_F}
 =a_{n,m}\sqrt{S_{n,m}},
\]
and dividing $\mathcal E_{n,m}$ by this quantity gives
\eqref{eq:energy-anova-t}.

\smallskip
\noindent\emph{Nesting of the two regressions.}
Consider a generic fitted value from the restricted model.  Define
\[
 \overline\gamma_X=\frac1n\sum_{i\le n}\gamma_i,
 \qquad
 \overline\gamma_Y=\frac1m\sum_{i>n}\gamma_i,
\]
and write
\[
 d_i=
 \begin{cases}
  \gamma_i-\overline\gamma_X,&i\le n,\\
  \gamma_i-\overline\gamma_Y,&i>n.
 \end{cases}
 \qquad
 \sum_{i\le n}d_i=\sum_{i>n}d_i=0.
\]
The restricted fitted values on $XX$, $XY$, and $YY$ edges are, respectively,
\[
 2\overline\gamma_X+d_i+d_j,
 \qquad
 \overline\gamma_X+\overline\gamma_Y+\theta+d_i+d_j,
 \qquad
 2\overline\gamma_Y+d_i+d_j.
\]
The fully interacted model reproduces these values by taking
$\alpha_i=\beta_i=d_i$,
$\eta=\overline\gamma_X+\overline\gamma_Y$, and
$\delta=\overline\gamma_X-\overline\gamma_Y$, while leaving $\theta$
unchanged.  These choices satisfy all four zero-sum constraints.  Thus the
restricted fitted space is contained in the fully interacted fitted space.

The restricted fitted rank is $N+1$, as shown above, while the full fitted
rank is $2N-1$.  The fully interacted model therefore adds $N-2$ fitted
directions, and
\[
 \nu_R-\nu_F=N-2.
\]
Let $\widehat f_R$ and $\widehat f_F$ denote the two fitted edge vectors.
Nested least squares gives the orthogonal decomposition
\[
 a_E-\widehat f_R
 =(a_E-\widehat f_F)+(\widehat f_F-\widehat f_R).
\]
Therefore
\[
 \mathrm{RSS}_R
 =\mathrm{RSS}_F+\mathrm{SS}_{\mathrm{int}},
 \qquad
 \mathrm{SS}_{\mathrm{int}}
 :=\norm{\widehat f_F-\widehat f_R}^2\ge0.
\]
The second term is the additional variation explained by allowing endpoint
effects to differ between within- and between-group pairs.  This proves
\eqref{eq:nested-rss-decomposition}.

Together, these calculations prove
Propositions~\ref{prop:energy-group-slope} and~\ref{prop:energy-anova-t}, as
well as the restricted-model formulas
\eqref{eq:restricted-two-sample-rss}--\eqref{eq:restricted-energy-t}.


\section*{AI-use statement}
During the preparation of this manuscript, the author used GPT-5.6 Sol for
language editing and exploratory assistance with proof
development.  The author independently checked every definition, statement,
and proof and takes full responsibility for the content of the manuscript.


\begin{thebibliography}{99}
\small

\bibitem[Bloznelis and G\"otze(2001)]{BloznelisGotze2001}
Bloznelis, M. and G\"otze, F. (2001).
Orthogonal decomposition of finite population statistics and its applications
to distributional asymptotics.
\emph{Annals of Statistics}, \textbf{29}, 899--917.
\url{https://doi.org/10.1214/aos/1009210694}.

\bibitem[Chakraborty and Zhang(2021)]{ChakrabortyZhang2021}
Chakraborty, S. and Zhang, X. (2021).
A new framework for distance and kernel-based metrics in high dimensions.
\emph{Electronic Journal of Statistics}, \textbf{15}, 5455--5522.
\url{https://doi.org/10.1214/21-EJS1889}.


 \bibitem[Dunkl(1978)]{Dunkl1978}
Dunkl, C. F. (1978).
An addition theorem for Hahn polynomials: The spherical functions.
\emph{SIAM Journal on Mathematical Analysis}, \textbf{9}, 627--637.
\url{https://doi.org/10.1137/0509043}.

\bibitem[Filmus(2016)]{Filmus2016}
Filmus, Y. (2016).
An orthogonal basis for functions over a slice of the Boolean hypercube.
\emph{Electronic Journal of Combinatorics}, \textbf{23}(1), P1.23.
\url{https://doi.org/10.37236/4567}.

\bibitem[Gao and Shao(2023)]{GaoShao2023}
Gao, H. and Shao, X. (2023).
Two sample testing in high dimension via maximum mean discrepancy.
\emph{Journal of Machine Learning Research}, \textbf{24}(304), 1--33.
\url{https://www.jmlr.org/papers/v24/22-1136.html}.

 \bibitem[Gottlieb(1966)]{Gottlieb1966}
Gottlieb, D. H. (1966).
A certain class of incidence matrices.
\emph{Proceedings of the American Mathematical Society}, \textbf{17},
1233--1237.
\url{https://doi.org/10.1090/S0002-9939-1966-0204305-9}.

\bibitem[Gretton et~al.(2005)Gretton, Bousquet, Smola, and Sch\"olkopf]{GrettonEtAl2005}
Gretton, A., Bousquet, O., Smola, A., and Sch\"olkopf, B. (2005).
Measuring statistical dependence with Hilbert--Schmidt norms.
In \emph{Algorithmic Learning Theory}, Lecture Notes in Computer Science,
\textbf{3734}, 63--77. Springer.
\url{https://doi.org/10.1007/11564089_7}.

\bibitem[Gretton et~al.(2007)Gretton, Fukumizu, Teo, Song, Sch\"olkopf, and Smola]{GrettonEtAl2008}
Gretton, A., Fukumizu, K., Teo, C. H., Song, L., Sch\"olkopf, B., and Smola,
A. J. (2007).
A kernel statistical test of independence.
In \emph{Advances in Neural Information Processing Systems 20 (NIPS 2007)}, 585--592.

\bibitem[Hoeffding(1948)]{Hoeffding1948}
Hoeffding, W. (1948).
A class of statistics with asymptotically normal distribution.
\emph{Annals of Mathematical Statistics}, \textbf{19}, 293--325.
 \url{https://doi.org/10.1214/aoms/1177730196}.


\bibitem[Huo and Sz\'ekely(2016)]{HuoSzekely2016}
Huo, X. and Sz\'ekely, G. J. (2016).
Fast computing for distance covariance.
\emph{Technometrics}, \textbf{58}, 435--447.
\url{https://doi.org/10.1080/00401706.2015.1054435}.

\bibitem[Kenny and La~Voie(1984)]{KennyLaVoie1984}
Kenny, D. A. and La Voie, L. (1984).
The social relations model.
\emph{Advances in Experimental Social Psychology}, \textbf{18}, 141--182.
\url{https://doi.org/10.1016/S0065-2601(08)60144-6}.

 \bibitem[Peccati(2004)]{Peccati2004}
Peccati, G. (2004).
Hoeffding--ANOVA decompositions for symmetric statistics of exchangeable
observations.
\emph{Annals of Probability}, \textbf{32}, 1796--1829.
\url{https://doi.org/10.1214/009117904000000405}.

 \bibitem[Sejdinovic et~al.(2013)Sejdinovic, Sriperumbudur, Gretton, and Fukumizu]{SejdinovicEtAl2013}
Sejdinovic, D., Sriperumbudur, B., Gretton, A., and Fukumizu, K. (2013).
Equivalence of distance-based and RKHS-based statistics in hypothesis testing.
\emph{Annals of Statistics}, \textbf{41}, 2263--2291.
\url{https://doi.org/10.1214/13-AOS1140}.

\bibitem[Serfling(1980)]{Serfling1980}
Serfling, R. J. (1980).
\emph{Approximation Theorems of Mathematical Statistics}.
Wiley, New York.

\bibitem[Shao and Zhang(2014)]{ShaoZhang2014}
Shao, X. and Zhang, J. (2014).
Martingale difference correlation and its use in high-dimensional variable
screening.
\emph{Journal of the American Statistical Association}, \textbf{109},
1302--1318.
\url{https://doi.org/10.1080/01621459.2014.887012}.

\bibitem[Song et~al.(2012)Song, Smola, Gretton, Bedo, and Borgwardt]{SongEtAl2012}
Song, L., Smola, A., Gretton, A., Bedo, J., and Borgwardt, K. (2012).
Feature selection via dependence maximization.
\emph{Journal of Machine Learning Research}, \textbf{13}, 1393--1434.

\bibitem[Szab\'o and Sriperumbudur(2018)]{SzaboSriperumbudur2018}
Szab\'o, Z. and Sriperumbudur, B. K. (2018).
Characteristic and universal tensor product kernels.
\emph{Journal of Machine Learning Research}, \textbf{18}(233), 1--29.
\url{https://jmlr.org/papers/v18/17-492.html}.

\bibitem[Sz\'ekely and Rizzo(2013)]{SzekelyRizzo2013}
Sz\'ekely, G. J. and Rizzo, M. L. (2013).
The distance correlation $t$-test of independence in high dimension.
\emph{Journal of Multivariate Analysis}, \textbf{117}, 193--213.
\url{https://doi.org/10.1016/j.jmva.2013.02.012}.

\bibitem[Sz\'ekely and Rizzo(2014)]{SzekelyRizzo2014}
Sz\'ekely, G. J. and Rizzo, M. L. (2014).
Partial distance correlation with methods for dissimilarities.
\emph{Annals of Statistics}, \textbf{42}, 2382--2412.
\url{https://doi.org/10.1214/14-AOS1255}.

\bibitem[Sz\'ekely et~al.(2007)Sz\'ekely, Rizzo, and Bakirov]{SzekelyRizzoBakirov2007}
Sz\'ekely, G. J., Rizzo, M. L., and Bakirov, N. K. (2007).
Measuring and testing dependence by correlation of distances.
\emph{Annals of Statistics}, \textbf{35}, 2769--2794.
\url{https://doi.org/10.1214/009053607000000505}.

\bibitem[Warner et~al.(1979)Warner, Kenny, and Stoto]{WarnerKennyStoto1979}
Warner, R. M., Kenny, D. A., and Stoto, M. (1979).
A new round robin analysis of variance for social interaction data.
\emph{Journal of Personality and Social Psychology}, \textbf{37}(10),
1742--1757.
\url{https://doi.org/10.1037/0022-3514.37.10.1742}.

\bibitem[Wang and Lindsay(2014)]{WangLindsay2014}
Wang, Q. and Lindsay, B. (2014).
Variance estimation of a general U-statistic with application to
cross-validation.
\emph{Statistica Sinica}, \textbf{24}, 1117--1141.
\url{https://doi.org/10.5705/ss.2012.215}.

 \bibitem[Zhang et~al.(2018)Zhang, Yao, and Shao]{ZhangYaoShao2018}
Zhang, X., Yao, S., and Shao, X. (2018).
Conditional mean and quantile dependence testing in high dimension.
\emph{Annals of Statistics}, \textbf{46}, 219--246.
\url{https://doi.org/10.1214/17-AOS1548}.

\bibitem[Zhu et~al.(2020)Zhu, Zhang, Yao, and Shao]{ZhuZhangYaoShao2020}
Zhu, C., Zhang, X., Yao, S., and Shao, X. (2020).
Distance-based and RKHS-based dependence metrics in high dimension.
\emph{Annals of Statistics}, \textbf{48}, 3366--3394.
\url{https://doi.org/10.1214/19-AOS1934}.

\end{thebibliography}

\end{document}
